% %
% manual.tex - Reference manual for gbemu
%                                       2026-10-02
%
%  Reference manual for the gbemu Game Boy / Game Boy Color emulator.
%  Layout follows the reference-manual style of the refart class: a narrow
%  text column with a wide left margin that carries headings and margin
%  notes.
%
\documentclass[twoside,a4paper]{refart}
\usepackage{makeidx}
\usepackage{ifthen}
\usepackage{longtable}
\usepackage{booktabs}
\usepackage{array}
\usepackage{fancyvrb}
\usepackage[hidelinks]{hyperref}

% ---------------------------------------------------------------------------
% Notation
% ---------------------------------------------------------------------------

% backslash in \tt font
\def\bs{\char'134 }

\newcommand{\ie}{i.\,e.,}
\newcommand{\eg}{e.\,g.,}

% Long identifiers must wrap: the text column is narrow and names like
% GB_HOTKEY_SAVE_STATE are wider than it.  \cs and \fname therefore insert a
% zero-width break opportunity between every token.
\usepackage{seqsplit}
\renewcommand*\seqinsert{\allowbreak}
\ExplSyntaxOn
\NewDocumentCommand{\brk}{m}{\tl_map_inline:nn {#1} {##1\allowbreak}}
\ExplSyntaxOff

% \cs{name}   -> "\name" in typewriter (macros without underscores)
% \fname{x}   -> "x" in typewriter, underscores safe (identifiers)
\DeclareRobustCommand\cs[1]{\texttt{\char`\\\brk{#1}}}
\DeclareRobustCommand\fname[1]{%
  \texttt{\edef\gbemudetok{\detokenize{#1}}%
         \expandafter\seqsplit\expandafter{\gbemudetok}}}

% Verbatim environments. "Code" is framed, "Trace" is not.
\DefineVerbatimEnvironment{Code}{Verbatim}{%
  fontsize=\footnotesize,
  frame=single,
  framesep=4pt,
  xleftmargin=2pt,
  baselinestretch=0.92
}
\DefineVerbatimEnvironment{Trace}{Verbatim}{%
  fontsize=\footnotesize,
  baselinestretch=0.92
}

% Column types for long tables
\newcolumntype{L}[1]{>{\raggedright\arraybackslash}p{#1}}
\newcolumntype{C}[1]{>{\centering\arraybackslash}p{#1}}
\renewcommand{\arraystretch}{1.05}
\setlength{\LTleft}{0pt}
\setlength{\LTright}{\fill}
% The refart text column is only about 316pt wide, so the default 6pt
% intercolumn glue pushes the wider tables past the right margin.
\setlength{\tabcolsep}{3pt}

% A boxed paragraph used for "gotchas"
\newsavebox{\warnbox}
\newenvironment{warning}{%
  \par\medskip\noindent
  \begin{lrbox}{\warnbox}%
  \begin{minipage}{\linewidth}%
  \hrule height 0.6pt
  \vspace{3pt}%
  \textbf{Note.}\hspace{0.5em}\ignorespaces
}{%
  \par\vspace{3pt}\hrule height 0.6pt
  \end{minipage}%
  \end{lrbox}%
  \noindent\box{\warnbox}\par\medskip
}

% Run-in notes.  refart's \attention and \seealso are margin commands that
% require an argument; the notes here are whole paragraphs of prose, so they
% are rendered in the text column under their own run-in label instead.
\renewcommand*\attention{\par\medskip\noindent{\bfseries !\enspace}\ignorespaces}
\renewcommand*\seealso{\par\medskip\noindent{\bfseries $\rightarrow$\enspace}\ignorespaces}

\title{gbemu\\
\large A Game Boy Emulator in C\\
\large Reference Manual}
\author{Konstantinos Despoinidis (KDesp73)\\
gbemu 0.1.0\\
MIT licensed\thanks{The full text of the MIT license ships with the
source tree in the file \fname{LICENSE}.}}

\date{October 2026}
\emergencystretch1em

\pagestyle{myfootings}
\markboth{gbemu: A Game Boy Emulator in C}
         {gbemu Reference Manual}

\makeindex

\setcounter{tocdepth}{2}

\begin{document}

\maketitle

\begin{abstract}
    \fname{gbemu} is a Game Boy and Game Boy Color emulator written in
    portable~C99.  The emulator core is a frontend-agnostic library: it
    knows nothing about windows, canvases or sound cards.  Everything the
    core cannot do by itself --- presenting a frame, reading a game pad,
    pushing audio samples --- is delegated to a small vtable that a host
    program fills in.  The same core therefore drives an SDL3 desktop
    build, a terminal front end that draws with ANSI half-blocks, a
    headless build used by the test harness, a WebAssembly build that runs
    in a browser, and an ESP32 firmware image.

    This manual documents the whole system: how to build it, how to run it,
    how the memory map, the cartridge mappers, the timer, the PPU and the
    APU are implemented, how save files and save states are laid out on
    disk, how to embed the library in your own program, and how the test
    suites are driven.  Where the implementation knowingly deviates from
    hardware, the deviation is stated rather than hidden.
\end{abstract}

\tableofcontents

\newpage

% ===========================================================================
\section{Introduction}
\label{intro}

\subsection{What gbemu is}

\index{gbemu}
\fname{gbemu} is a from-scratch emulator for the original monochrome Game
Boy and, partially, for the Game Boy Color.  It is written in C99 with no
mandatory third-party dependency: the core library talks to nothing but
the C standard library.  Optional back ends add SDL3 (desktop), emscripten
(browser) or ESP-IDF (microcontroller) support, but each one is a
separate, replaceable translation unit.

The design goal is \emph{testability first}.  Every hardware subsystem is
a small, independently callable state machine that advances by an explicit
number of cycles, so that conformance ROMs can be run at full speed with
no window and no sound card.  The second goal is \emph{transparency}: the
source is small enough to read in an afternoon, and this manual explains
every register, every timing constant and every deliberate deviation from
real hardware.

\marginlabel{Scope:}
The emulator implements the DMG feature set, the four supported memory bank
controllers, the four audio channels, and the interrupt system.  Game Boy
Color support is limited to the CPU speed switch; the CGB-specific
palette, video and DMA registers are not implemented and are listed in
section~\ref{cgb-limits}.

\subsection{Design goals}
\label{goals}

\begin{description}

\item[Frontend independence]\index{frontend!independence}
    The core exposes a four-function vtable (\fname{init},
    \fname{render}, \fname{poll_events}, \fname{destroy}) through
    \fname{gb_frontend}, plus an optional hotkey callback.  A frontend
    never pokes at emulator internals; it receives a frame buffer and a
    memory bus.  This is what makes the headless test build possible.

\item[Cycle-driven hardware]\index{cycle-driven}
    The CPU does not "run" the hardware; hardware is advanced explicitly.
    \fname{gb_cpu_step} reports how many T-cycles an instruction took
    and calls \fname{gb_machine_tick} at each machine-cycle boundary, so
    the timer, the PPU, the APU and the OAM DMA engine all observe bus
    activity at the right slot in time.

\item[Generated opcode tables]\index{code generation}
    Opcode mnemonics and cycle counts come from a machine-readable table
    and are regenerated into \fname{src/opcodes.h} by
    \fname{scripts/generate_opcodes.py}.  Hand-maintained opcode lists
    drift; generated ones do not.

\item[Documented from the source]\index{documentation}
    The public header \fname{src/gbemu.h} carries structured comments
    (\fname{//@func}, \fname{//@param}, \fname{//@desc}, \ldots) that are
    turned into browsable API documentation by \fname{make docs}.  This
    manual covers the parts a header comment cannot: why the code looks the
    way it does, and what it does not do.

\item[Portability]\index{portability}
    Only the core library has to be portable.  The timing helpers
    (\fname{gb_get_time_ns}, \fname{gb_sleep_ns}) switch between
    POSIX and FreeRTOS behind \fname{\#ifdef}, and the ESP32 build compiles
    the core freestanding, with no ESP-IDF headers.
\end{description}

\subsection{Feature summary}
\label{features}

\begin{longtable}{@{}L{0.22\textwidth}L{0.72\textwidth}@{}}
\toprule
\textbf{Area} & \textbf{What is implemented}\\
\midrule
\endhead
CPU &
  Full SM83 instruction set: all 256 base opcodes and all 256
  \fname{CB}-prefixed opcodes, with hardware-accurate cycle counts,
  the \fname{EI} one-instruction delay, the \fname{HALT} bug and
  interrupt dispatch.\\
Memory bus &
  Complete address decoding including echo RAM, the unusable region,
  I/O register routing, the joypad matrix, OAM DMA and a flat I/O
  shadow array for registers that are not modelled.\\
Cartridges &
  ROM-only, MBC1 (including MBC1M multicarts), MBC2, MBC3 (with an RTC
  register stub) and MBC5.  Battery saves for every mapper that has
  battery-backed RAM.\\
Timer &
  \fname{DIV}/\fname{TIMA}/\fname{TMA}/\fname{TAC} with a 16-bit internal
  counter, a bit-selected frequency clock and a falling-edge detector,
  including the delayed \fname{TIMA} reload and the \fname{DIV}-write
  trigger.\\
PPU &
  Background, window and sprite rendering to a 160$\times$144 ARGB frame
  buffer, scanline mode transitions (OAM, transfer, H-blank, V-blank),
  \fname{LY}/\fname{LYC} coincidence and the four STAT interrupt sources.\\
Joypad &
  \fname{P1} at \fname{0xFF00} with the two select lines, active-low
  button state and joypad interrupt generation.\\
OAM DMA &
  \fname{0xFF46} transfers copy 160 bytes at one byte per machine cycle
  after a two-cycle start delay, with restart queueing and OAM blocking
  during the copy.\\
Audio &
  Four-channel synthesis (two squares with duty, a wave channel reading
  channel~3's sample RAM, and a 15/7-bit LFSR noise channel), NR51/NR50
  panning and master volume, naive decimation to the output rate and a
  single-producer/single-consumer ring buffer.\\
Serial &
  Output is intercepted: a write of \fname{0x81} to \fname{0xFF02}
  prints the pending byte, and writes to cartridge RAM above
  \fname{0xA004} print printable bytes.  This is what makes the
  Blargg, mooneye and SameSuite suites reportable.\\
Persistence &
  Battery saves (\fname{<rom>.sav}) and full-machine save states
  (\fname{<rom>.state}) with version and ROM-hash validation.\\
Frontends &
  SDL3, terminal (ANSI truecolour half-blocks), headless, WebAssembly
  and ESP32.\\
\bottomrule
\end{longtable}

\subsection{Accuracy and known limitations}
\label{limits-intro}

An emulator that documents only its successes is not much use to anyone
trying to decide whether a given test failure is the emulator's fault.
The known deviations are therefore listed where they belong, in the
chapter of the subsystem they affect:

\begin{itemize}
\item \textbf{CPU:} \fname{STOP} consumes its operand byte and performs the
      CGB speed switch, but it does not stop the CPU, the divider or the
      APU.\index{STOP}
\item \textbf{Bus:} unmapped I/O reads return \fname{0x00} rather than
      \fname{0xFF}; VRAM and OAM are not blocked during PPU modes~2 and~3;
      the OAM DMA source page is not masked or validated; \fname{0xFF46}
      does not read back the value written to it.
\item \textbf{Timer:} writing \fname{TAC} does not generate the
      falling edge that real hardware produces.
\item \textbf{PPU:} scanlines are rendered in one step at dot~504 rather
      than pixel by pixel, mode~3 has a fixed length, \fname{LCDC} bit~0 is
      ignored, the sprite priority flag (OAM attribute bit~7) is parsed
      but not applied, scanline~0 is not drawn after the LCD is switched
      on, and there is no OAM-scan corruption bug.
\item \textbf{APU:} the frame sequencer period, the noise channel shift
      and the wave-channel volume register are all wrong in ways
      section~\ref{apu-limits} quantifies; the volume envelope is not
      wired up.
\item \textbf{CGB:} no palette RAM, no VRAM banking, no GDMA/HDMA, no
      boot ROM.  See section~\ref{cgb-limits}.
\end{itemize}

\subsection{Repository layout}
\label{repo-layout}

\begin{longtable}{@{}L{0.24\textwidth}L{0.70\textwidth}@{}}
\toprule
\textbf{Path} & \textbf{Contents}\\
\midrule
\endhead
\fname{src/} &
  The emulator core.  \fname{gbemu.h} is the entire public API;
  \fname{frontend.h} declares the vtable.  \fname{cpu.c},
  \fname{instr.c}, \fname{opcodes.h} (generated), \fname{bus.c},
  \fname{timer.c}, \fname{ppu.c}, \fname{apu.c}, \fname{inter.c},
  \fname{loop.c}, \fname{memops.c} and \fname{save.c} make up the
  implementation.  Nothing in this directory includes SDL or any
  windowing header.\\
\fname{frontend/} &
  One file per front end: \fname{sdl.c}, \fname{terminal.c},
  \fname{headless.c}, \fname{wasm.c}, \fname{esp32.c}.  Exactly one of
  these is compiled into a given host program; the build system picks
  the default.\\
\fname{apps/} &
  Host programs.  \fname{cli.c} is the desktop/CLI entry point,
  \fname{wasm.c} is the emscripten entry point with the exported
  functions the web shell calls.\\
\fname{make/} &
  GNU make fragments: \fname{common.mk} (compiler and flags),
  \fname{libs.mk} (artifacts), \fname{utils.mk} (housekeeping),
  \fname{wasm.mk}, \fname{esp32.mk}, \fname{docs.mk}.\\
\fname{scripts/} &
  \fname{generate_opcodes.py} regenerates \fname{src/opcodes.h};
  \fname{test} is the test-suite driver.\\
\fname{tests/} &
  Checked-out test ROMs: \fname{gb-test-roms} (Bla'rg),
  \fname{mooneye-test-suite} and \fname{SameSuite}.\\
\fname{docs/} &
  This manual (\fname{manual.tex}), developer notes under
  \fname{docs/dev/}, and the generated API reference
  (\fname{tiny.docs.json} plus \fname{index.html}).\\
\fname{web/} &
  \fname{shell.html}, the emscripten shell template that becomes
  \fname{emu-wasm.html}.\\
\bottomrule
\end{longtable}

\subsection{Conventions used in this manual}
\label{conventions}

\begin{description}

\item[Identifiers] are set in typewriter font: \fname{gb_bus_write}.
\item[Hexadecimal constants] are written \texttt{0xFF40} with an
      \texttt{0x} prefix and four digits for addresses; masks and
      immediate values may be shorter.
\item[Time bases] are named explicitly.  A \emph{dot} (a system cycle) is
      a \fname{GB_CPU_FREQ} unit at 8.388608\,MHz.  A \emph{T-cycle} is
      a CPU clock unit, equal to two dots at normal speed.  An
      \emph{M-cycle} (machine cycle) is four T-cycles, equal to one bus
      access.  Section~\ref{time-bases} tabulates the conversions.
\item[Addresses] are given as full I/O addresses where the I/O space is
      involved (\texttt{0xFF05}) and as cartridge-window offsets where the
      mapper is involved (\texttt{\$2000--\$3FFF}), matching the
      conventions of the Game Boy hardware documentation.
\item[\texttt{DONE}/\texttt{TODO}] items in the source \fname{TODO} file
      are quoted when they explain a missing feature.
\end{description}

\subsection{License and authorship}

\index{license}
The project is released under the MIT license.  The public header
declares this too, through the \fname{//@license} documentation marker.
Contributions are welcome; the fastest way to make progress is to pick an
item from the \fname{TODO} file or a failing conformance test and send a
patch.

\clearpage

% ===========================================================================
\section{Building gbemu}
\label{building}

\subsection{Prerequisites}

The core needs only a C99 compiler.  The rest depends on the target:

\begin{longtable}{@{}L{0.26\textwidth}L{0.68\textwidth}@{}}
\toprule
\textbf{Target} & \textbf{Requirements}\\
\midrule
\endhead
Core library, CLI &
  \textbf{gcc} (or any C99 compiler).  On macOS with Homebrew this is
  the system compiler.  \textbf{pkg-config} is used to locate SDL3.\\
SDL3 front end &
  SDL3 development files, discovered with
  \cs{pkg-config --cflags --libs sdl3}.  The makefiles additionally point
  \cs{PKG\_CONFIG\_PATH} at \fname{/opt/homebrew/lib/pkgconfig}, which is
  where Homebrew installs it on Apple Silicon.\\
Terminal / headless &
  Nothing beyond libc; useful on systems without a display server.\\
WebAssembly &
  \textbf{emcc} from the Emscripten SDK.  Only the core, the WASM front
  end and \fname{apps/wasm.c} are compiled; SDL is not involved.\\
ESP32 &
  The Espressif cross toolchain, on \cs{PATH} after sourcing
  \cs{\$HOME/esp/esp-idf/export.sh}.  The core cross-compiles
  freestanding; \fname{frontend/esp32.c} additionally needs a full
  ESP-IDF component environment.\\
Documentation &
  \textbf{pdflatex} for this manual, \textbf{tinydocs-cli} for the API
  reference, \textbf{rgbds} (for SameSuite) and \textbf{wla-dx} (for
  mooneye) if you want to rebuild the test ROMs.\\
\bottomrule
\end{longtable}

\subsection{Build targets}
\label{make-targets}

The top-level \fname{Makefile} includes five fragments and defines
\cs{help} as the default goal.  Every target below carries a \cs{\#\#}
comment so that \cs{make help} lists it:

\begin{longtable}{@{}L{0.26\textwidth}L{0.68\textwidth}@{}}
\toprule
\textbf{Target} & \textbf{Description}\\
\midrule
\endhead
\cs{all} &
  The default build: checks the toolchain, creates \fname{build/}, then
  builds the static library, the shared library and the CLI binary.
  Equivalent to \cs{check\_tools static shared gbemu-cli}.\\
\cs{static} &
  Archives the core objects into \fname{libgbemu.a} with
  \cs{ar rcs}.\\
\cs{shared} &
  Links the same objects into \fname{libgbemu.so} with \cs{-shared}.\\
\cs{gbemu-cli} &
  Links \fname{apps/cli.c} together with the selected front end against
  \fname{libgbemu.a}.  This is the program you actually run.\\
\cs{check\_tools} &
  Verifies that \textbf{gcc} exists (fatal otherwise) and warns if
  \textbf{bear} is missing, since that only affects the
  \cs{compile\_commands.json} target.\\
\cs{clean} &
  Removes \fname{build/}, the binary and both libraries.\\
\cs{distclean} &
  \cs{clean} plus editor backups, core dumps and
  \fname{dist/*.tar.gz}.\\
\cs{dist} &
  Packs \fname{src}, \fname{frontend}, \fname{apps}, the
  \fname{Makefile} and \fname{README.md} into
  \fname{dist/gbemu-cli.tar.gz}.\\
\cs{compile\_commands.json} &
  Regenerates the compile database used by clangd and IDEs, via
  \cs{bear -- make all}.\\
\cs{docs} &
  Runs \fname{tinydocs-cli} over \fname{src/gbemu.h} and
  \fname{src/frontend.h} and writes \fname{docs/tiny.docs.json} plus an
  HTML page.\\
\cs{manual} &
  Compiles this document: three \textbf{pdflatex} passes around a
  \textbf{makeindex} run, with the output collected in
  \fname{docs/manual/} (see section~\ref{manual-build}).\\
\cs{manual.clean} &
  Removes \fname{docs/manual/}.\\
\cs{install} / \cs{uninstall} &
  Copy, respectively remove, the binary in \cs{PREFIX}, which defaults
  to \cs{/usr/local/bin}.  Both need write access, i.e.\ usually
  \cs{sudo}.\\
\cs{wasm} &
  Cross-compiles the core plus the WASM front end with \textbf{emcc}
  into \fname{web/emu-wasm.html} and its sidecar files.\\
\cs{wasm.serve} &
  Rebuilds the WASM target and serves \fname{web/} on
  \cs{http://localhost:8080} with \cs{python3 -m http.server}.\\
\cs{wasm.clean} &
  Removes the WASM objects and every \fname{emu-wasm.*} artefact.\\
\cs{esp32-lib} &
  Cross-compiles the core for the microcontroller into
  \fname{build/esp32/libgbemu-<target>.a}.\\
\cs{esp32-clean} &
  Removes \fname{build/esp32}.\\
\bottomrule
\end{longtable}

\marginlabel{Variables:}
The build accepts three meaningful variables.  \cs{type=RELEASE} switches
to an optimised, sanitizer-free build.  \cs{TERM=1} and \cs{HEADLESS=1}
select the terminal and headless front ends instead of SDL3.
\cs{WASM\_ROM=<path>} preloads a ROM into the WASM virtual file system,
and \cs{ESP32\_TARGET=<chip>} picks the cross toolchain.  \cs{PREFIX}
overrides the install directory.

\subsection{Build types}

\cs{make/common.mk} treats \cs{type=RELEASE} specially:

\begin{description}

\item[Debug (default)] \cs{-DDEBUG -ggdb} plus
  \cs{-fsanitize=address,undefined} in both the compile and the link
  line.  This is the build you want while developing: the sanitizers
  catch out-of-bounds accesses in the frame buffer and in the DMA code
  that would otherwise show up as a corrupted picture two frames later.

\item[Release] \cs{make all type=RELEASE} compiles with \cs{-O3} and no
  sanitizers.  Everything always compiles with \cs{-Wall}, \cs{-Isrc} and
  \cs{-fPIC} (the latter so the objects can go into the shared library).
\end{description}

\subsection{Choosing a front end at build time}

\cs{make/common.mk} inspects two variables and sets exactly one front-end
source file, plus the matching \cs{-D} macro that \fname{apps/cli.c}
tests:

\begin{longtable}{@{}C{0.16\textwidth}L{0.22\textwidth}L{0.54\textwidth}@{}}
\toprule
\textbf{Variable} & \textbf{Define} & \textbf{Front end}\\
\midrule
\endhead
(not set) & --- & \fname{frontend/sdl.c}, compiled together with the
  SDL3 flags reported by \cs{pkg-config}\\
\cs{TERM=1} & \cs{-DEMU\_TERM} & \fname{frontend/terminal.c}; no SDL
  dependency at all, so this is the build to use on a headless box\\
\cs{HEADLESS=1} & \cs{-DEMU\_HEADLESS} & \fname{frontend/headless.c};
  initialises nothing, renders nothing, never terminates on its own.  The
  test harness uses this.\\
\bottomrule
\end{longtable}

The front-end factory function differs per build, which is why
\fname{apps/cli.c} declares three of them and links only one:

\begin{Code}
gb_frontend* frontend_sdl_create(gb_apu* apu);
gb_frontend* frontend_headless_create(void);
gb_frontend* frontend_terminal_create(void);
\end{Code}

Note that the SDL constructor takes a \fname{gb_apu*} because its audio
callback needs the ring buffer; the other two take nothing and produce no
sound.

\subsection{Build outputs}

\cs{make all} leaves three artefacts in the repository root:

\begin{longtable}{@{}L{0.24\textwidth}L{0.70\textwidth}@{}}
\toprule
\textbf{Artefact} & \textbf{Description}\\
\midrule
\endhead
\fname{gbemu-cli} &
  The executable.  Takes a ROM path as its only argument.\\
\fname{libgbemu.a} &
  Static library containing the core only --- no front end, no
  \cs{main}.  This is the artefact to link into an application of your
  own.\\
\fname{libgbemu.so} &
  The same objects as a shared library, for run-time loading or for
  embedding in another language.\\
\bottomrule
\end{longtable}

\subsection{The WebAssembly target}
\label{wasm-build}

\cs{make wasm} compiles every core source \emph{except} \fname{loop.c}
(the browser supplies its own frame loop), adds \fname{frontend/wasm.c}
and \fname{apps/wasm.c}, and links with \textbf{emcc}.  The link options
matter:

\begin{Code}
-s ALLOW_MEMORY_GROWTH
-s INITIAL_MEMORY=16777216
-s EXPORTED_RUNTIME_METHODS='["FS","HEAPU8"]'
-s EXPORTED_FUNCTIONS='["_main","_wasm_set_key","_wasm_load_rom","_wasm_is_running"]'
--shell-file web/shell.html
\end{Code}

Two consequences follow.  \fname{FS} is required because the page writes
the chosen ROM into the virtual file system as \fname{/rom.gb} before
calling \fname{_wasm_load_rom}.  \fname{HEAPU8} is exported because the
front end wraps a pointer into the WASM linear memory to hand it to
JavaScript.  \cs{INITIAL\_MEMORY} of 16\,MiB leaves room for a large ROM
plus its bank-aligned allocation; \cs{ALLOW\_MEMORY\_GROWTH} lets a huge
cartridge still fit.

Debug builds add \cs{-s ASSERTIONS=2} so that a null pointer or a
misaligned access traps instead of corrupting memory silently.

Passing \cs{WASM\_ROM=path/to/rom.gb} adds emcc's
\cs{-\,-preload-file} option with the target
\fname{@/rom.gb}, so the page boots straight into a game without a file
picker.

\subsection{The ESP32 target}
\label{esp32-build}

\cs{make esp32-lib} compiles the core for the microcontroller and
archives it:

\begin{Code}
make esp32-lib                    # xtensa-esp32-elf-gcc
make ESP32_TARGET=esp32s3 esp32-lib
make ESP32_TARGET=esp32c3 esp32-lib   # riscv32-esp-elf-gcc
\end{Code}

Supported \cs{ESP32\_TARGET} values are \cs{esp32}, \cs{esp32s2},
\cs{esp32s3} and \cs{esp32c3}; anything else is a hard make error.  The
compile flags are deliberately minimal: \cs{-O2}, \cs{\-ffunction-sections}
and \cs{\-fdata-sections} so the linker can garbage-collect unused code,
plus \cs{-DESP\_PLATFORM}, which is the marker that switches the portable
paths inside \fname{loop.c} and \fname{bus.c}:

\begin{itemize}
\item \fname{gb_sleep_ns} uses \cs{vTaskDelay} instead of
      \cs{nanosleep}, selected by a compile-time probe for
      \cs{<freertos/FreeRTOS.h>}.
\item The \cs{EMU\_NOSLEEP} environment variable is never read; frame
      pacing is unconditional.
\item The Bla'rg trace file side effect is compiled out.
\end{itemize}

The core needs no ESP-IDF headers, so the archive links into any IDF
project.  \fname{frontend/esp32.c} is \emph{not} part of this target: it
needs the full component environment, and section~\ref{esp32-frontend}
describes how to use it from your own firmware.

\subsection{Building this manual}
\label{manual-build}

\cs{make manual} compiles \fname{docs/manual.tex} into
\fname{docs/manual/}; the target is defined in \fname{make/docs.mk} and
shells out to the toolchain below.  The document uses only packages from
a standard TeX Live installation (\cs{refart} with \cs{longtable},
\cs{booktabs}, \cs{array}, \cs{fancyvrb} and \cs{hyperref}), and it has to
run the LaTeX pass more than once so that the index and the page
references settle:

\begin{Code}
cd docs
pdflatex -interaction=nonstopmode -output-directory=manual manual.tex
makeindex -o manual/manual.idx manual/manual.idx
pdflatex -interaction=nonstopmode -output-directory=manual manual.tex
pdflatex -interaction=nonstopmode -output-directory=manual manual.tex
\end{Code}

\cs{make manual} does exactly this, including the \cs{-output-directory},
so the \fname{.aux}, \fname{.log} and \fname{.idx} files never appear next
to the source.  If you invoke \cs{pdflatex} by hand without that flag,
compile in \fname{docs/} instead and clean up the intermediate files
afterwards with \fname{latexmk -c} or by hand.

\attention
The index is only complete if \cs{makeindex} runs \emph{after} the first
LaTeX pass has written \fname{manual.idx} and \emph{before} the last one.
Skipping it produces a document that compiles cleanly and prints an empty
index.

The layout is the reference-manual style: the text occupies a narrow
column on the right, and the wide left margin carries \cs{marginlabel}
annotations, \cs{seealso} cross references and \cs{attention} notes.  If
you extend the document, that is where the asides belong.

\clearpage

% ===========================================================================
\section{Running the emulator}
\label{running}

\subsection{Command line}

The CLI takes exactly one argument, the path to a ROM file:

\begin{Code}
$ ./gbemu-cli roms/"Tetris (World) (Rev 1).gb"
\end{Code}

With no argument it prints a usage message and exits with status~1:

\begin{Code}
Please provide a rom
Usage: ./gbemu-cli <ROM>
\end{Code}

Paths are used verbatim, including spaces --- quote them.  The ROM's
directory matters, because both the battery save and the save state are
written next to the ROM (section~\ref{persistence}).

\subsection{Environment variables}

\begin{longtable}{@{}L{0.20\textwidth}L{0.13\textwidth}L{0.59\textwidth}@{}}
\toprule
\textbf{Variable} & \textbf{Default} & \textbf{Meaning}\\
\midrule
\endhead
\cs{EMU\_SCALE} & \cs{4} &
  Integer window scale for the SDL front end; the window is
  \cs{160*scale} by \cs{144*scale}.  Values below 1 are clamped to 1.
  \cs{EMU\_SCALE=3} gives a 480$\times$432 window.\\
\cs{EMU\_NOSLEEP} & unset &
  If the variable exists at all --- its value is ignored --- the
  emulator does not sleep between frames and runs as fast as it can.
  This is what the test harness sets, and what you want when
  measuring emulation speed.  It is not read at all in the ESP32
  build.\\
\bottomrule
\end{longtable}

\subsection{Controls}
\label{controls}

The joypad state is a pair of bytes on the bus, \fname{joypad_dpad} and
\fname{joypad_buttons}, both \emph{active low}: a set bit means released.
Only the low nibble of each byte is used.

\begin{longtable}{@{}C{0.14\textwidth}C{0.22\textwidth}C{0.28\textwidth}C{0.28\textwidth}@{}}
\toprule
\textbf{Button} & \textbf{Field} & \textbf{Bit} & \textbf{SDL3 key}\\
\midrule
\endhead
Right & \fname{joypad_dpad} & \texttt{0x01} & \cs{SDLK\_RIGHT}\\
Left  & \fname{joypad_dpad} & \texttt{0x02} & \cs{SDLK\_LEFT}\\
Up    & \fname{joypad_dpad} & \texttt{0x04} & \cs{SDLK\_UP}\\
Down  & \fname{joypad_dpad} & \texttt{0x08} & \cs{SDLK\_DOWN}\\
\midrule
A     & \fname{joypad_buttons} & \texttt{0x01} & \cs{SDLK\_Z}\\
B     & \fname{joypad_buttons} & \texttt{0x02} & \cs{SDLK\_X}\\
Select & \fname{joypad_buttons} & \texttt{0x04} & \cs{SDLK\_BACKSPACE}\\
Start & \fname{joypad_buttons} & \texttt{0x08} & \cs{SDLK\_RETURN}\\
\bottomrule
\end{longtable}

Any SDL key press raises the joypad interrupt flag; the ESP32 front end
detects the 1-to-0 transitions properly instead, and the terminal front
end raises it once per poll unconditionally (a known quirk,
section~\ref{terminal-frontend}).

\subsection{Hotkeys}

The front end never decides what a hotkey \emph{means}.  It recognises the
input and emits an action through the \fname{on_hotkey} callback; the
application decides that a save action should call \fname{gb_save_state}.

\begin{longtable}{@{}C{0.24\textwidth}L{0.34\textwidth}L{0.34\textwidth}@{}}
\toprule
\textbf{Action} & \textbf{SDL3} & \textbf{ESP32}\\
\midrule
\endhead
\fname{GB_HOTKEY_SAVE_STATE} & \cs{SDLK\_F5} & \cs{START + SELECT}
  held for about half a second\\
\fname{GB_HOTKEY_LOAD_STATE} & \cs{SDLK\_F9} & \cs{START + SELECT}
  held for about two and a half seconds\\
\bottomrule
\end{longtable}

Terminal, headless and WASM builds install no hotkey callback at all, so
the keys do nothing there.  The full wiring is shown in
section~\ref{embedding-hotkeys}.

\subsection{Battery saves and save states}

Both kinds of file are derived from the ROM path by appending an extension:

\begin{longtable}{@{}L{0.24\textwidth}L{0.34\textwidth}L{0.34\textwidth}@{}}
\toprule
\textbf{File} & \textbf{Written} & \textbf{Contains}\\
\midrule
\endhead
\fname{<rom>.sav} &
  Automatically when \fname{gb_loop} returns, for cartridges whose
  header declares a battery &
  Cartridge RAM, and the RTC registers for MBC3 timer carts\\
\fname{<rom>.state} &
  On \cs{F5} &
  The complete machine state: CPU, bus, timer, PPU, APU\\
\bottomrule
\end{longtable}

On start-up the emulator loads \fname{<rom>.sav} if it exists.  Save states
are loaded only when you press the load hotkey, and are rejected unless
they were written by the same emulator version from the same ROM
(section~\ref{state-validation}).

\subsection{Shutting down}

The emulation loop keeps a local \fname{running} flag which only the front
end can clear, through the \cs{running} out-parameter of
\fname{poll_events}.  The SDL front end clears it on \cs{SDL\_EVENT\_QUIT}
or a window-close request; the terminal front end clears it on
\cs{Ctrl-C}, on a bare escape, or on \cs{q}.  Once the loop exits,
\cs{main} writes the battery save, then destroys the front end and returns.

Two consequences are worth knowing.  Input is only sampled once per frame,
so quitting takes effect within roughly one frame.  And because the headless
front end never clears the flag, a headless process runs until it is
killed --- which is exactly what the test harness relies on.

\clearpage

% ===========================================================================
\section{Architecture}
\label{architecture}

\subsection{Layers}

The system is three layers deep, and each layer only knows the one below
it:

\begin{Code}
  +---------------------------------------------------------------+
  |  host program          apps/cli.c      apps/wasm.c           |
  |  main(), state objects, hotkey policy, exit handling          |
  +-------------------------------+-------------------------------+
                                  |  gb_frontend* (vtable)
  +-------------------------------v-------------------------------+
  |  front end                  frontend/sdl.c, terminal.c,      |
  |  pixels, audio, input        headless.c, wasm.c, esp32.c      |
  +-------------------------------+-------------------------------+
                                  |  gb_bus*, gb_apu*
  +-------------------------------v-------------------------------+
  |  core library               src/cpu.c instr.c bus.c timer.c  |
  |  (libgbemu.a / .so)         src/ppu.c apu.c inter.c loop.c    |
  |                              src/memops.c save.c              |
  +---------------------------------------------------------------+
\end{Code}

The core never calls into a front end except through the four vtable
slots, and a front end never calls into the core except
\fname{gb_apu_buf_pop} and the bus read/write functions it needs for
input.  That is the whole contract.

\subsection{Module map}
\label{module-map}

\begin{longtable}{@{}L{0.20\textwidth}L{0.74\textwidth}@{}}
\toprule
\textbf{File} & \textbf{Responsibility}\\
\midrule
\endhead
\fname{cpu.c} &
  Register initialisation, register dumping, the flag helpers, the
  register-index helpers, and \fname{gb_cpu_step} --- the fetch,
  decode and dispatch entry point, including the \fname{EI} delay and
  the HALT bug.\\
\fname{instr.c} &
  The instruction set.  Eleven handler functions cover the base
  opcodes; one more covers the \fname{CB} prefix.  All of them are
  \cs{static}; the public surface is \fname{gb_instr}.\\
\fname{opcodes.h} &
  Generated.  Two enumerations of opcode names and two
  \cs{static inline} functions that return the cycle count for a
  given opcode and condition.\\
\fname{bus.c} &
  Address decoding for the whole 16-bit space, I/O register routing,
  ROM loading and header parsing, the four memory bank controllers,
  the joypad matrix, serial output and the OAM DMA engine.  It also
  owns \fname{gb_machine_tick}, the function that advances every
  other component.\\
\fname{timer.c} &
  The divider and the timer, plus the two portability helpers
  \fname{gb_get_time_ns} and \fname{gb_sleep_ns}.\\
\fname{ppu.c} &
  The mode state machine, the scanline renderer and the video
  register file.\\
\fname{apu.c} &
  The frame sequencer, the four channels, the mixer and the output
  ring buffer.\\
\fname{inter.c} &
  \fname{gb_handle_inter_interrupts}: turns hardware request flags
  into IF bits, wakes a halted CPU and performs the five-machine-cycle
  dispatch sequence.\\
\fname{loop.c} &
  \fname{gb_loop}: the frame-budgeted main loop, the render handoff,
  input polling and frame pacing.\\
\fname{memops.c} &
  Two helpers, \fname{gb_fetch8} and \fname{gb_fetch16}, which tick
  one machine cycle per byte read and assemble little-endian 16-bit
  values.\\
\fname{save.c} &
  Battery saves and save states.\\
\bottomrule
\end{longtable}

\subsection{The front-end vtable}
\label{vtable}

\fname{src/frontend.h} declares the whole interface between core and
host:

\begin{Code}
typedef enum {
    GB_HOTKEY_SAVE_STATE, // Quick-save the full machine state (F5)
    GB_HOTKEY_LOAD_STATE, // Restore the last quick-save (F9)
} gb_hotkey;

struct gb_frontend {
    void* priv;

    bool (*init)(gb_frontend* fe, int width, int height);
    void (*render)(gb_frontend* fe, const uint32_t* buffer, int width, int height);
    void (*poll_events)(gb_frontend* fe, gb_bus* bus, bool* running);
    void (*destroy)(gb_frontend* fe);

    void* hotkey_ctx;
    void (*on_hotkey)(void* ctx, gb_hotkey key);
};
\end{Code}

The contract, call by call:

\begin{longtable}{@{}L{0.18\textwidth}L{0.32\textwidth}L{0.42\textwidth}@{}}
\toprule
\textbf{Slot} & \textbf{Called} & \textbf{Contract}\\
\midrule
\endhead
\fname{init} &
  Once, by \cs{main}, with \cs{160} and \cs{144} &
  Returns \cs{true} on success.  Allocate everything private here;
  return \cs{false} to abort start-up, which makes \cs{main} exit
  with status~1.\\
\fname{render} &
  Once per frame, only when \cs{ppu->frame\_ready} was set &
  \cs{buffer} points at \cs{160*144} ARGB words in row-major order.
  The core clears the flag after the call, so the front end must
  copy what it needs rather than keep the pointer.\\
\fname{poll_events} &
  Once per frame, after rendering &
  Drain the platform event queue, write the joypad bits into
  \cs{bus->joypad\_dpad} / \cs{bus->joypad\_buttons} (active low),
  raise \cs{bus->joypad\_interrupt} when a button was newly pressed,
  and clear \cs{*running} to end the session.\\
\fname{destroy} &
  Once, by \cs{main} after the loop &
  Release everything \cs{init} allocated.  The core does not touch
  the vtable afterwards.\\
\fname{on_hotkey} &
  From \cs{poll\_events}, if non-NULL &
  The front end signals an intent; the application performs the
  save or load.  May be \fname{NULL}.\\
\bottomrule
\end{longtable}

\marginlabel{Ordering:}
Rendering strictly precedes event polling, and both happen after the
cycle budget for the frame is spent.  Input therefore takes effect on the
following frame, at most 16.7\,ms later.

\subsection{The execution model}
\label{exec-model}

The core is a cycle-driven state machine.  Nothing in it is
event-driven: there are no callbacks from the timer into the CPU, no
interrupts-as-preemption, no coroutines.  Each component is a plain struct
plus a \cs{step} function that takes a cycle count.

The call graph of one frame of \fname{gb_loop}:

\begin{Code}
gb_loop(cpu, bus, timer, ppu, apu, fe)
 |
 +- gb_get_time_ns()                       frame start timestamp
 |
 +- while (frame_cycles < GB_CYCLES_PER_FRAME):
 |    |
 |    +- gb_cpu_step(cpu, bus)  --------------------------> T-cycles
 |    |     +- promote ime_scheduled -> ime
 |    |     +- if halted: gb_machine_tick(bus,4); return 4
 |    |     +- gb_machine_tick(bus, 4)                    M-cycle 1
 |    |     +- opcode = gb_bus_read(bus, pc); pc++       (skipped on halt bug)
 |    |     +- gb_instr(cpu, bus, opcode)
 |    |           +- 0xCB ? instr_cb() : handler chain
 |    |           +- handler calls gb_machine_tick(bus,4) per M-cycle
 |    |           +- handlers call gb_fetch8/16 (which tick too)
 |    +- gb_handle_interrupts(cpu, bus, ppu, timer) -----> T-cycles
 |    |     +- latch PPU/timer/joypad requests into IF (0xFF0F)
 |    |     +- wake the CPU if IF & IE != 0
 |    |     +- if IME: push PC, jump to 0x40 + 8*bit, cost 20 T-cycles
 |    +- frame_cycles += t_cycles * (double_speed ? 1 : 2)
 |
 +- if (ppu->frame_ready): fe->render(...); frame_ready = false
 +- fe->poll_events(fe, bus, &running)
 +- gb_get_time_ns(); gb_sleep_ns(remaining of 1/60 s) unless EMU_NOSLEEP
\end{Code}

Four properties of this design are worth stating explicitly, because they
explain most of the code.

\begin{description}

\item[Cycles are returned, never accumulated]\index{cycle accounting}
    \fname{gb_cpu_step} returns the number of T-cycles an instruction
    consumed, taken from the generated table \emph{before} execution
    starts.  Handlers report only whether they handled the opcode.
    Because the returned total is fixed, conditional branches are
    charged their taken or untaken cost based on a
    \emph{pre-execution} evaluation of the flag --- which is why
    \fname{get_condition_met} is called before any handler runs.

\item[Hardware advances at machine-cycle boundaries]
    \fname{gb_machine_tick} is called once for the opcode fetch and
    once for every subsequent bus access.  It steps the DMA engine, the
    timer, the PPU and the APU.  A load from \fname{[HL]} costs a
    machine cycle because the register-index helper calls it
    explicitly, which is what makes the \cs{BIT n,(HL)} timing come
    out right.

\item[Interrupts are checked between instructions]
    \fname{gb_handle_interrupts} runs after every completed
    instruction, so worst-case interrupt latency is one instruction.
    Only one interrupt is dispatched per call; the remaining pending
    bits stay in IF for the next instruction boundary.

\item[A frame is a cycle budget, not a V-blank]
    The inner loop runs until \cs{GB\_CYCLES\_PER\_FRAME} dots have
    been spent, and renders if --- and only if --- the PPU raised
    \cs{frame\_ready} during that budget.  If a frame overruns, the
    next render simply shows the newer picture; the flag is cleared
    on read, so a second V-blank inside one budget is dropped rather
    than queued.
\end{description}

\subsection{Time bases}
\label{time-bases}

\begin{longtable}{@{}L{0.22\textwidth}L{0.20\textwidth}L{0.52\textwidth}@{}}
\toprule
\textbf{Unit} & \textbf{Magnitude} & \textbf{Notes}\\
\midrule
\endhead
Dot (system cycle) & 8.388608\,MHz &
  \cs{GB\_CPU\_FREQ}.  The PPU and the APU count dots; the frame
  budget is expressed in dots, which is why it is speed-mode
  independent.\\
T-cycle & 4.194304\,MHz &
  Two dots at normal speed, one dot in CGB double-speed mode.  The CPU
  and the timer live in this domain; the timer is
  \emph{not} rescaled, so \fname{DIV} ticks twice as fast in real time
  in double-speed mode, exactly as on hardware.\\
M-cycle & 4 T-cycles &
  One bus access.  \fname{gb_machine_tick} is always called with
  \cs{4}.\\
Frame & \cs{GB\_CYCLES\_PER\_FRAME} =
  \cs{GB\_CPU\_FREQ/60} = 139\,810 dots &
  The inner-loop budget.\\
Frame wall clock & \cs{GB\_FRAME\_TIME\_NS} =
  $10^9/60$ = 16\,666\,666\,ns &
  The pacing target.\\
Scanline & 912 dots = 456 T-cycles &
  See section~\ref{ppu-modes}.\\
Interrupt dispatch & 20 T-cycles = 5 M-cycles &
  See section~\ref{interrupts}.\\
\bottomrule
\end{longtable}

\subsection{State objects and ownership}
\label{ownership}

An application owns five objects and keeps them alive for the whole
session:

\begin{Code}
gb_cpu   cpu;    // registers, IME, halt state
gb_bus   bus;    // memory map, cartridge, MBC state, joypad, DMA
gb_timer timer;  // DIV/TIMA/TMA/TAC and the internal counter
gb_ppu   ppu;    // video registers and the frame buffer
gb_apu   apu;    // sound registers, channel state, ring buffer
\end{Code}

They must be wired together before the first tick:

\begin{Code}
bus.timer = &timer;
bus.ppu   = &ppu;
bus.apu   = &apu;
\end{Code}

\attention
\fname{gb_machine_tick} guards each of these with an
\cs{if (bus->ppu)} style test.  If you forget a pointer, the machine
does not crash --- it silently stops advancing that component, and you
get a black screen with no error message.

Two other ownership rules:

\begin{itemize}
\item \fname{gb_bus_load_rom} allocates \cs{bus->rom} and
      \cs{bus->sram} itself and frees any previous allocation, so a
      reload is safe.  Nothing else allocates.
\item The PPU frame buffer is the only large allocation in the state
      objects (\cs{160*144*4} bytes, about 90\,KiB).  It is
      deliberately placed last in the struct so that save states can
      serialise the register prefix with a single
      \cs{offsetof}.
\end{itemize}

\subsection{Symbol prefixing and \cs{GBEMU\_STRIP\_PREFIX}}
\label{prefix}

Every public symbol carries a \fname{gb_} prefix and every macro a
\fname{GB_} prefix: \fname{gb_bus_write}, \fname{gb_ppu_step},
\fname{GB_SCREEN_WIDTH}.  The enum constants
(\fname{GB_FLAG_Z}, \fname{GB_PPU_MODE_VBLANK}) follow the same
rule.  This keeps the header usable in projects that already have a
symbol called \fname{loop} or \fname{version}.

If you are writing a single-file program that links nothing else, you can
ask the header to undo the prefixing:

\begin{Code}
cc -DGBEMU_STRIP_PREFIX my_game.c -I src libgbemu.a -o my_game -lSDL3
\end{Code}

After that, \fname{cpu_init} expands to \fname{gb_cpu_init},
\fname{CPU} to \fname{gb_cpu}, \fname{SCREEN_WIDTH} to
\fname{GB_SCREEN_WIDTH}, and so on.  The mechanism is a flat list of
object-like \cs{\#define} aliases placed at the very end of
\fname{gbemu.h}, after the real declarations, so the ABI never changes.

\attention
This deliberately gives up namespacing.  It redefines very common words
(\fname{loop}, \fname{version}, \fname{STR}, \fname{loop}) process-wide.
Pass it on the command line as shown rather than defining it in a source
file, so that it cannot leak into unrelated translation units, and never
enable it in a project that also uses another library with those names.

Only the prefixed spellings are aliased for the types that are also
declared in \fname{frontend.h}; including either header alone gives you
the same short names.

\clearpage

% ===========================================================================
\section{The CPU}
\label{cpu}

\subsection{Register file}

\fname{gb_cpu} mirrors the hardware with four anonymous unions, one per
register pair:

\begin{Code}
typedef struct {
    union { struct { uint8_t f; uint8_t a; }; uint16_t af; };
    union { struct { uint8_t c; uint8_t b; }; uint16_t bc; };
    union { struct { uint8_t e; uint8_t d; }; uint16_t de; };
    union { struct { uint8_t l; uint8_t h; }; uint16_t hl; };
    uint16_t sp;
    uint16_t pc;
    bool ime;             // Interrupt Master Enable
    bool ime_scheduled;   // EI has a 1-instruction delay
    bool halted;
    bool halt_bug;        // HALT with IME=0 and an interrupt pending
} gb_cpu;
\end{Code}

\marginlabel{Byte order:}
In every pair the \emph{low} register occupies the low byte, so
\fname{af = 0x11B0} means \fname{f == 0xB0} and \fname{a == 0x11}.  The
16-bit value therefore matches the hardware's \fname{AF} while the
storage order is reversed; code that needs \fname{A} from \fname{af}
uses \cs{cpu->af >> 8}.  \fname{POP AF} masks the low nibble of
\fname{F}, as required.

The post-boot register set:

\begin{longtable}{@{}L{0.22\textwidth}L{0.20\textwidth}L{0.52\textwidth}@{}}
\toprule
\textbf{Register} & \textbf{Value} & \textbf{Rationale}\\
\midrule
\endhead
\cs{AF} & \cs{0x11B0} &
  The Game Boy Color boot ROM leaves \fname{A = 0x11}, which selects
  the enhanced register set.  The application overrides it to
  \cs{0x01} for DMG-only carts by checking header byte
  \fname{0x0143}.\\
\cs{BC} & \cs{0x0013} & Boot ROM value\\
\cs{DE} & \cs{0x00D8} & Boot ROM value\\
\cs{HL} & \cs{0x014D} & Boot ROM value\\
\cs{SP} & \cs{0xFFFE} & Just below the I/O block, as after boot\\
\cs{PC} & \cs{0x0100} &
  Execution begins at the cartridge entry point; no boot ROM is
  emulated.\\
\midrule
\cs{timer.internal\_counter} & \cs{0xABCC} &
  Chosen so that the upper byte, which is what \fname{DIV} reads,
  equals the post-boot \cs{0xAB}.\\
\cs{bus.io[0x0F]} & \cs{0x01} &
  \fname{IF}: a V-blank request left over from the last scanline of
  the boot ROM.  Set by the application, not by the core.\\
\cs{bus.ie} & \cs{0x00} &
  \fname{IE}: no interrupt source is enabled out of boot.\\
\cs{joypad\_buttons}, \cs{joypad\_dpad} & \cs{0x0F} &
  All buttons released, given the active-low encoding.\\
\bottomrule
\end{longtable}

The first six rows are set by \fname{gb_cpu_init} and
\fname{gb_timer_init}; the rest are the host program's
responsibility, and both \fname{apps/cli.c} and \fname{apps/wasm.c} set
them identically.

\subsection{Flags}

\begin{longtable}{@{}C{0.18\textwidth}C{0.18\textwidth}L{0.56\textwidth}@{}}
\toprule
\textbf{Flag} & \textbf{Bit} & \textbf{Set when}\\
\midrule
\endhead
\cs{GB\_FLAG\_Z} & 7 & The result of the operation is zero\\
\cs{GB\_FLAG\_N} & 6 & The last operation was a subtraction\\
\cs{GB\_FLAG\_H} & 5 & Half-carry or half-borrow out of bit~3\\
\cs{GB\_FLAG\_C} & 4 & Carry or borrow out of bit~7\\
\bottomrule
\end{longtable}

Bits~0--3 of \fname{F} do not exist and always read as zero;
\fname{POP AF} enforces this.  \fname{gb_flag_set} and
\fname{gb_flag_get} are the only sanctioned way to touch them, and all
of \fname{instr.c} goes through them --- 137 call sites, which makes the
flag semantics auditable in one place.

\subsection{Register-index helpers}
\label{regindex}

Most opcodes encode their operand in three bits.  \fname{gb_get_reg_by_index}
and \fname{gb_set_reg_by_index} turn that index into a register:

\begin{longtable}{@{}C{0.10\textwidth}L{0.28\textwidth}L{0.52\textwidth}@{}}
\toprule
\textbf{Index} & \textbf{Register} & \textbf{Note}\\
\midrule
\endhead
\cs{0} & \cs{B} & \\
\cs{1} & \cs{C} & \\
\cs{2} & \cs{D} & \\
\cs{3} & \cs{E} & \\
\cs{4} & \cs{H} & \\
\cs{5} & \cs{L} & \\
\cs{6} & \cs{(HL)} &
  The only index with a side effect: a read or write of the bus at
  address \fname{HL}, costing one machine cycle.\\
\cs{7} & \cs{A} & \\
\bottomrule
\end{longtable}

Because index~6 pays a machine cycle, \cs{BIT n,(HL)} costs 12 T-cycles
(three machine cycles: fetch, prefix byte, memory read) and
\cs{RES}/\cs{SET n,(HL)} cost 16 (four), while the register forms cost
one and two respectively.  The timing falls out of the data path rather
than being special-cased.

\subsection{Instruction dispatch}

\fname{gb_instr} is the dispatcher.  It looks the cycle cost up first,
then walks a table of eleven handlers until one of them claims the
opcode:

\begin{Code}
static instrfn instr_handlers[] = {
    instr_8bit_arithm, instr_16bit_arithm, instr_bit_flag,
    instr_bit_shift,   instr_bitwise,      instr_carry,
    instr_interrupts,  instr_jumps,        instr_load,
    instr_misc,        instr_stack
};

int gb_instr(gb_cpu* cpu, gb_bus* bus, uint8_t opcode) {
    if (opcode == OP_PREFIX) return instr_cb(cpu, bus);
    int cycles = get_opcode_cycles(opcode, get_condition_met(cpu, opcode));
    for (size_t i = 0; i < N; ++i)
        if (instr_handlers[i](cpu, bus, opcode)) return cycles;
    return 0;  // no handler claimed it
}
\end{Code}

Each handler has the same shape: one \cs{switch} over the opcode, with
fall-through \cs{case} labels grouping the register variants of one
mnemonic, and a \cs{default: return false;} to decline.  Inside a body the
register index is recovered arithmetically with \cs{opcode \& 0x07} (source
operand) and \cs{(opcode >> 3) \& 0x07} (destination, bit position, or
\fname{RST} vector class).

\begin{longtable}{@{}L{0.28\textwidth}C{0.16\textwidth}L{0.46\textwidth}@{}}
\toprule
\textbf{Handler} & \textbf{Opcodes} & \textbf{Contents}\\
\midrule
\endhead
\fname{instr_8bit_arithm} & 61 &
  \cs{ADD ADC SUB SBC AND XOR OR CP} against every register and
  immediate, plus \cs{INC} / \cs{DEC} in all their forms\\
\fname{instr_load} & 92 &
  Every \cs{LD} variant: register, immediate, \cs{[HL]},
  \cs{[HL+]}, \cs{[HL-]}, \cs{LDH}, \cs{[a16]}, \cs{(C)}, and the
  high-memory forms\\
\fname{instr_16bit_arithm} & 13 &
  \cs{ADD HL,r16}, \cs{ADD SP,e8}, \cs{INC r16}, \cs{DEC r16}\\
\fname{instr_bitwise} & 28 &
  \cs{AND OR XOR CPL} against every register and immediate, including
  \cs{CP} against \cs{(HL)}\\
\fname{instr_bit_flag} & 0 &
  An empty placeholder: a \cs{switch} with only a \cs{default} that
  declines.  It sits in the dispatch table, so every instruction pays
  one extra indirect call for it.\\
\fname{instr_bit_shift} & 4 &
  Only the accumulator forms \cs{RLCA}, \cs{RRCA}, \cs{RLA}, \cs{RRA};
  the register forms are reached through \cs{CB}\\
\fname{instr_jumps} & 30 &
  \cs{JP}, \cs{JR}, \cs{CALL}, \cs{RET}, \cs{RETI}, \cs{RST} in all
  their conditional and unconditional forms\\
\fname{instr_carry} & 2 & \cs{CCF}, \cs{SCF}\\
\fname{instr_interrupts} & 3 & \cs{DI}, \cs{EI}, \cs{HALT}\\
\fname{instr_stack} & 8 &
  \cs{PUSH BC DE HL AF} and \cs{POP BC DE HL AF}; \cs{LD [n16],SP},
  \cs{ADD SP,e8} and \cs{LD HL,SP+e8} live in \fname{instr_load} and
  \fname{instr_16bit_arithm} instead\\
\fname{instr_misc} & 3 & \cs{DAA}, \cs{NOP}, \cs{STOP}\\
\midrule
\textbf{Base opcodes} & \textbf{244} &
  256 minus the \cs{CB} prefix byte and the eleven illegal opcodes.
  The remaining twelve reach the fatal path in \fname{gb_cpu_step}.\\
\bottomrule
\end{longtable}

If no handler claims an opcode, \fname{gb_instr} returns zero and
\fname{gb_cpu_step} treats that as fatal:

\begin{Code}
fprintf(stderr, "Unhandled opcode: 0x%02X at PC: 0x%04X\n", opcode, pc - 1);
exit(1);
\end{Code}

That path is reachable only through the eleven illegal opcodes
(\cs{0xD3}, \cs{0xDB}, \cs{0xDD}, \cs{0xE3}, \cs{0xE4}, \cs{0xEB},
\cs{0xEC}, \cs{0xED}, \cs{0xF4}, \cs{0xFC}, \cs{0xFD}), which the opcode
table names and prices but no handler implements.

\subsection{Generated opcode tables}
\label{opcode-tables}

\fname{src/opcodes.h} is machine-generated by
\fname{scripts/generate_opcodes.py} from a JSON opcode table and must
never be edited by hand.  It contains:

\begin{itemize}
\item An anonymous enumeration of all 256 base opcodes, named
      \cs{OP\_<mnemonic>\_<operands>}.
\item An anonymous enumeration of all 256 \fname{CB}-prefixed opcodes,
      with the prefix byte folded into the value
      (\cs{OP\_CB\_RLC\_B = 0xCB00}).
\item \cs{get\_opcode\_cycles(opcode, condition\_met)} and
      \cs{get\_cb\_opcode\_cycles(cb\_opcode)}, both \cs{static inline}
      switches returning the T-cycle cost.
\end{itemize}

Naming rules, worth knowing when reading \fname{instr.c}:

\begin{longtable}{@{}L{0.30\textwidth}L{0.60\textwidth}@{}}
\toprule
\textbf{Rule} & \textbf{Example}\\
\midrule
\endhead
Mnemonic in upper case, operands appended &
  \cs{OP\_ADD\_A\_B} for \cs{ADD A,B}\\
\cs{\$} becomes \cs{S} & \cs{OP\_RST\_S38} for \cs{RST \$38}\\
Non-immediate operands gain \cs{\_IND} & \cs{OP\_LD\_HL\_IND\_B} for \cs{LD (HL),B}\\
Post-increment operands gain \cs{\_INC} & \cs{OP\_LD\_HL\_INC\_A} for \cs{LD (HL+),A}\\
Post-decrement operands gain \cs{\_DEC} & \cs{OP\_LD\_A\_HL\_DEC} for \cs{LD A,(HL-)}\\
Size hints stay verbatim & \cs{OP\_JP\_a16}, \cs{OP\_ADD\_SP\_e8},
  \cs{OP\_STOP\_n8}\\
\bottomrule
\end{longtable}

The cycle functions encode the taken/untaken distinction in the table
itself.  A conditional branch emits
\cs{return condition\_met ? 12 : 8;}, and the caller must therefore
evaluate the condition \emph{before} running the handler --- flags are
clobbered by the time the handler returns.

The generated \cs{OP\_*} constants exist mainly for readability;
\fname{instr.c} matches on raw opcode numbers and recovers operand indices
with shifts and masks, which is both faster and easier to verify against
the reference tables.

\subsection{\cs{CB}-prefixed instructions}

The 256 \fname{CB} opcodes are \emph{not} dispatched through a table.
\fname{instr_cb} decodes the second byte as bit fields:

\begin{longtable}{@{}C{0.20\textwidth}C{0.20\textwidth}L{0.50\textwidth}@{}}
\toprule
\textbf{Field} & \textbf{Bits} & \textbf{Meaning}\\
\midrule
\endhead
\fname{reg_idx} & \cs{0x07} & As in section~\ref{regindex}\\
\fname{bit_pos} & \cs{(b>>3) \& 0x07} & Bit number for \cs{BIT},
  \cs{RES}, \cs{SET}, and the shift amount for the rotate group\\
\fname{group} & \cs{(b>>6) \& 0x03} & Instruction family\\
\bottomrule
\end{longtable}

\begin{longtable}{@{}C{0.16\textwidth}L{0.34\textwidth}L{0.40\textwidth}@{}}
\toprule
\textbf{Group} & \textbf{Instructions} & \textbf{Notes}\\
\midrule
\endhead
\cs{0x00} & \cs{RLC}, \cs{RRC}, \cs{RL}, \cs{RR}, \cs{SLA}, \cs{SRA},
  \cs{SWAP}, \cs{SRL} & Chosen by \cs{bit\_pos}; always writes back\\
\cs{0x01} & \cs{BIT n,r} & Writes nothing; sets \fname{Z} from the bit,
  clears \fname{N}, sets \fname{H}; returns immediately\\
\cs{0x02} & \cs{RES n,r} & Clears one bit, writes back\\
\cs{0x03} & \cs{SET n,r} & Sets one bit, writes back\\
\bottomrule
\end{longtable}

The operand is fetched \emph{unconditionally}, before the group is
dispatched, which is precisely what makes the \cs{[HL]} forms cost more:
12 T-cycles for \cs{BIT n,(HL)} and 16 for \cs{RES}/\cs{SET n,(HL)},
against 8 for the register forms.

\subsection{Instruction timing}
\label{instr-timing}

All instruction costs are multiples of four T-cycles:

\begin{longtable}{@{}C{0.18\textwidth}L{0.74\textwidth}@{}}
\toprule
\textbf{T-cycles} & \textbf{Representative instructions}\\
\midrule
\endhead
4  & \cs{NOP}, \cs{SCF}, \cs{CCF}, \cs{DAA}, \cs{CPL}, \cs{HALT}, \cs{STOP},
     \cs{RLCA}, all single-byte arithmetic on registers
     (for example \cs{INC B}, \cs{ADD A,B}), and every
     \cs{JR d8}\\
8  & Loads and arithmetic with an immediate (\cs{LD A,d8}, \cs{ADD A,d8},
     \cs{CP d8}), \cs{RET}, \cs{RET cc} untaken, \cs{LD [HLI],A},
     \cs{LD A,[HLD]}, \cs{LD (C),A}, \cs{LD SP,HL}, and every
     \cs{CB} operation on a register\\
12 & \cs{JR cc}, \cs{POP}, \cs{LDH}, \cs{LDH (a8),A}, \cs{LD HL,SP+e8},
     \cs{JP cc} untaken, \cs{CALL cc} untaken, \cs{BIT n,(HL)}\\
16 & \cs{JP a16}, \cs{JP cc} taken, \cs{CALL cc} taken, \cs{RET cc} taken,
     \cs{PUSH}, \cs{ADD SP,e8}, \cs{LD r8,[HL]}, \cs{LD [HL],r8},
     \cs{RES}/\cs{SET n,(HL)}\\
20 & \cs{LD [n16],SP}\\
24 & \cs{CALL a16}\\
\bottomrule
\end{longtable}

\marginlabel{Two prices per conditional:}
The table entry for a conditional instruction is written
\cs{return condition\_met ? taken : untaken;}.  Of the 240 non-\cs{CB}
opcodes, 145 cost 4, 60 cost 8, 18 cost 16, 15 cost 12, and the
multi-byte instructions \cs{CALL a16} and \cs{LD [n16],SP} cost 24 and 20.
Of the 256 \cs{CB} opcodes, 224 cost 8, 24 cost 16 and 8 cost 12.

\marginlabel{Invariant:}
For every opcode, the number of \cs{gb\_machine\_tick(bus, 4)} calls made
along the execution path (opcode fetch plus each operand or memory access)
equals the table's T-cycle count divided by four.  Nothing checks this at
run time, so an added access in a handler must be paid for by removing
one elsewhere.

\subsection{Interrupts}
\label{interrupts}

Hardware raises requests as booleans; \fname{gb_handle_interrupts} latches
them into \fname{IF}, decides whether to wake or to dispatch, and returns
the T-cycles consumed.

\begin{longtable}{@{}C{0.08\textwidth}C{0.18\textwidth}L{0.24\textwidth}L{0.16\textwidth}C{0.24\textwidth}@{}}
\toprule
\textbf{IF} & \textbf{Source} & \textbf{Raised by} & \textbf{Vector} &
\textbf{Priority}\\
\midrule
\endhead
\cs{0} & V-blank & \fname{ppu->vblank_interrupt} & \cs{0x0040} & 1 (highest)\\
\cs{1} & LCD STAT & \fname{ppu->stat_interrupt} & \cs{0x0048} & 2\\
\cs{2} & Timer & \fname{timer->interrupt_requested} & \cs{0x0050} & 3\\
\cs{3} & Serial & \emph{nothing} & \cs{0x0058} & 4\\
\cs{4} & Joypad & \fname{bus->joypad_interrupt} & \cs{0x0060} & 5\\
\bottomrule
\end{longtable}

The serial bit is never set by the hardware model: serial output is
handled by interception (section~\ref{serial}), not by a shift
register.  Priority is the lowest set bit, and exactly one interrupt is
dispatched per call.

The register semantics:

\begin{itemize}
\item \fname{IF} at \texttt{0xFF0F} reads back as \cs{io[0x0F] | 0xE0}:
      bits~5--7 read as one.  Writes are masked with \cs{0x1F}.
\item \fname{IE} at \texttt{0xFFFF} is plain storage in all eight bits,
      read and written unmasked.
\item The pending mask is \cs{IF \& IE \& 0x1F}.  The trailing \cs{0x1F} is
      essential: without it the unused \fname{IF} bits would appear to be
      permanently pending.
\item A pending interrupt clears \cs{halted} whether or not \fname{IME}
      is set.
\end{itemize}

When \fname{IME} is set and something is pending, dispatch clears \cs{IF}
and \fname{IME} immediately and then consumes five machine cycles
(20 T-cycles) in this order:

\begin{enumerate}
\item two internal machine cycles (no bus access)
\item decrement \cs{SP} by two, one machine cycle, after the two
      internal cycles
\item write the \emph{high} byte of \fname{PC} to \cs{SP+1}
\item write the \emph{low} byte of \fname{PC} to \cs{SP}
\item set \cs{PC} to \cs{0x0040 + 8*bit} --- one final internal machine
      cycle
\end{enumerate}

The two stack writes go through \fname{gb_bus_write}, so a game with
its stack page in cartridge RAM sees the real bus behaviour, including
mapper side effects.

\attention
The interrupt-enable tests are the strictest part of the timing model,
because they observe \cs{IE} and the joypad register mid-frame.  The
matching conformance tests are run by \fname{scripts/test}; the current
failure list is in the \fname{TODO} file.

\subsection{\cs{EI}, \cs{HALT}, the HALT bug and \cs{STOP}}
\label{halt-ei}

Two booleans model the enable sequence:

\begin{description}

\item[\cs{ime\_scheduled}] is set by \cs{EI} and cleared by \cs{DI}.
\item[\cs{ime}] is set from \cs{ime\_scheduled} at the \emph{top} of
      \fname{gb_cpu_step}, before the halt check.
\end{description}

This gives \cs{EI} its one-instruction delay: the instruction following
\fname{EI} executes with interrupts still masked, and the first
instruction after \emph{that} can be interrupted.
\cs{DI} clears both flags, so \cs{EI} immediately followed by \cs{DI}
never enables interrupts.  \cs{RETI} sets \fname{IME} immediately, as on
hardware, and does not touch \cs{ime\_scheduled}.

\fname{HALT} has three distinct outcomes:

\begin{longtable}{@{}L{0.34\textwidth}L{0.56\textwidth}@{}}
\toprule
\textbf{Condition} & \textbf{Behaviour}\\
\midrule
\endhead
Nothing pending & The CPU sleeps; each \fname{gb_cpu_step} burns one
  machine cycle and returns~4 without touching \fname{PC}.\\
Pending and \cs{IME=1} & The CPU does not sleep.  \fname{PC} stays at
  \cs{HALT+1} and the pending interrupt is dispatched normally after
  the instruction.\\
Pending and \cs{IME=0} & \emph{The HALT bug.}  The CPU does not sleep and
  \fname{halt_bug} is set.\\
\bottomrule
\end{longtable}

The HALT bug is modelled faithfully: the next opcode fetch reads its byte
without incrementing \fname{PC}, so that byte is read twice and every
subsequent operand read is shifted back one position.  The flag is
one-shot --- it is cleared by the fetch that consumes it --- which is
exactly the hardware behaviour for the one-byte \cs{HALT}.

\cs{STOP} (\cs{0x10}) consumes its operand byte, then flips
\fname{bus->double_speed} if \fname{KEY1} bit~0 was set, which is the
only way the CGB changes CPU speed.  It does \emph{not} stop the CPU,
halt the divider or mute the APU; those remain open items in the
\fname{TODO} list.

\marginlabel{Open item:}
The interaction between the \cs{EI} delay and \cs{HALT} is the one
interrupt test that is known to fail (SameSuite
\fname{interrupt/ei_delay_halt}), and is tracked in the
\fname{TODO} file.

\clearpage

% ===========================================================================
\section{The memory bus}
\label{bus}

\subsection{Address decoding}

\fname{gb_bus_read} and \fname{gb_bus_write} decode the full 16-bit
space.  The read path, in order:

\begin{longtable}{@{}L{0.24\textwidth}L{0.66\textwidth}@{}}
\toprule
\textbf{Range} & \textbf{Read behaviour}\\
\midrule
\endhead
\texttt{0x0000--0x7FFF} & Cartridge ROM, through the mapper's bank
  resolution (\fname{bus_read_rom})\\
\texttt{0x8000--0x9FFF} & \cs{bus->vram[addr - 0x8000]}\\
\texttt{0xA000--0xBFFF} & Cartridge RAM, through the mapper
  (\fname{bus_read_sram})\\
\texttt{0xC000--0xDFFF} & \cs{bus->wram[addr - 0xC000]}\\
\texttt{0xE000--0xFDFF} & Echo RAM: \cs{bus->wram[addr - 0xE000]}, the
  same 8\,KiB seen at \texttt{0xC000}\\
\texttt{0xFE00--0xFEA0} & Sprite attribute table, or \cs{0xFF} while an
  OAM DMA transfer is copying\\
\texttt{0xFF00--0xFF7F} & I/O dispatch (section~\ref{io-dispatch})\\
\texttt{0xFF80--0xFFFE} & \cs{bus->hram[addr - 0xFF80]}, 127 bytes\\
\texttt{0xFFFF} & \cs{bus->ie}\\
fallthrough & \cs{0xFF} --- this is exactly the unusable region
  \texttt{0xFEA0--0xFEFF}, the only gap in the 16-bit space\\
\bottomrule
\end{longtable}

The write path mirrors it, with two differences: writes to the ROM window
go to the mapper and never store anything, and writes to the unusable
region are dropped explicitly instead of falling through to a read
default.

\marginlabel{Not modelled:}
VRAM and OAM are \emph{not} blocked during PPU modes~2 and~3, and no
access restrictions are implemented for any other region either.  Real
games do rely on the restrictions in a handful of places; see
section~\ref{ppu-limits} for what that costs.

\subsection{I/O dispatch}
\label{io-dispatch}

The I/O window is not a flat array.  Reads and writes test specific
addresses and ranges first, then fall back to a 128-byte shadow array
\fname{bus->io[0x80]}:

\begin{longtable}{@{}L{0.26\textwidth}L{0.64\textwidth}@{}}
\toprule
\textbf{Address} & \textbf{Handler}\\
\midrule
\endhead
\texttt{0xFF00} & Joypad \fname{P1}\\
\texttt{0xFF04--0xFF07} & \fname{gb_timer_read} / \fname{gb_timer_write}\\
\texttt{0xFF10--0xFF3F} & \fname{gb_apu_read} / \fname{gb_apu_write}
  (including the wave RAM at \texttt{0xFF30}--\texttt{0xFF3F})\\
\texttt{0xFF40--0xFF4B} & \fname{gb_ppu_read} / \fname{gb_ppu_write},
  except that a write to \texttt{0xFF46} starts an OAM DMA and returns
  immediately\\
\texttt{0xFF4D} & \cs{KEY1}: bit~7 is synthesised from
  \cs{bus->double\_speed}, bit~0 is the speed-switch request\\
\texttt{0xFF0F} & \cs{IF}: reads OR \cs{0xE0}, writes mask \cs{0x1F}\\
otherwise & \cs{bus->io[addr - 0xFF00]}\\
\bottomrule
\end{longtable}

\marginlabel{Gotcha:}
Every one of the three component pointers is tested before use.  A bus
whose \cs{timer}, \cs{ppu} or \cs{apu} pointer is \fname{NULL} falls through
to the shadow array, so those registers appear to work --- they just do
nothing.  The same shadow array is what unmapped registers such as
\texttt{0xFF03}, \texttt{0xFF27} and the unimplemented CGB registers land
in; they read back whatever was last written, initially \cs{0x00} rather
than the \cs{0xFF} real hardware returns.

\subsection{Loading a ROM}
\label{rom-loading}

\fname{gb_bus_load_rom} opens the file, sizes it, allocates and clears
the ROM buffer, and inspects three header bytes:

\begin{Code}
bus->rom_banks = next power of two >= ceil(file_size / 0x4000), min 2;
bus->rom_size  = rom_banks * 0x4000;      // zero-padded
bus->rom       = calloc(1, bus->rom_size);
fread(bus->rom, 1, file_size, f);
\end{Code}

Deriving the bank count from the \emph{file size} rather than from header
byte \fname{0x0148} is a deliberate simplification: because the count is a
power of two, the mapper's usual masking \cs{bank \& (rom\_banks - 1)} is
always in range, whatever the header claims.  It also means a header that
lies about its size is ignored.

\begin{longtable}{@{}L{0.16\textwidth}L{0.30\textwidth}L{0.46\textwidth}@{}}
\toprule
\textbf{Offset} & \textbf{Meaning} & \textbf{Use}\\
\midrule
\endhead
\fname{0x0143} & CGB flag &
  \cs{0x80} (CGB enhanced) or \cs{0xC0} (CGB only) set
  \cs{apu->cgb}; the application also uses it to decide whether to
  force \cs{A = 0x01}.\\
\fname{0x0147} & Cartridge type &
  Selects the mapper (section~\ref{mappers})\\
\fname{0x0149} & RAM size &
  Sizes the cartridge RAM allocation (section~\ref{sram-sizing})\\
\midrule
\fname{0x0104} & Nintendo logo &
  Read \emph{only} for the MBC1M multicart heuristic: a 1\,MiB MBC1
  cart whose four 256\,KiB pages contain the logo in at least three of
  them is a multicart.\\
\bottomrule
\end{longtable}

Header byte \fname{0x0148} (ROM size), the checksums and the SGB flag are
not read at all.

\subsection{Sanitising the interrupt vectors}
\label{vector-sanitisation}

Conformance ROMs deliberately leave the five interrupt vectors empty or
pointing at unusable memory, then install their own handlers at run time.
The loader rewrites any vector that is still empty when the ROM is
loaded:

\begin{itemize}
\item a vector byte of \cs{0x00} becomes \cs{0xC9} (\cs{RET});
\item a vector starting with \cs{0xC3} (\cs{JP a16}) whose target lies in
      \fname{0x8000--0xBFFF} or \fname{0xFE00--0xFF7F} becomes
      \cs{0xC9}.
\end{itemize}

Only \cs{JP a16} is recognised, and the WRAM and HRAM targets are left
alone on purpose --- the Bla'rg suites copy their runtime into
\fname{0xC000} upwards and jump there.  Because the hash used for save
states is computed over the buffer \emph{after} this rewrite, states stay
consistent with each other.

\subsection{The joypad register}
\label{joypad}

\fname{P1} at \texttt{0xFF00} is the classic two-row matrix.  Writes
store only the two select bits:

\begin{longtable}{@{}L{0.26\textwidth}L{0.64\textwidth}@{}}
\toprule
\textbf{Bit} & \textbf{Meaning}\\
\midrule
\endhead
6, 7 & Read as one; ignored on write\\
5 (\cs{P15}) & Active low.  Clear selects the button row\\
4 (\cs{P14}) & Active low.  Clear selects the D-pad row\\
3--0 & The selected row, active low: a set bit is a released button\\
\bottomrule
\end{longtable}

Reads start from \cs{0xC0 | select} and then \emph{and} the selected
rows into the low nibble, so selecting both rows yields \cs{0x0F} and
selecting neither yields the intersection of both rows.

A write also raises the joypad interrupt when the newly selected row
contains a pressed button, which approximates the hardware's
``selected line goes low'' edge.

\attention
The comment in \fname{gbemu.h} describing bits~4--7 of the joypad bytes
is stale.  Only bits~0--3 of \fname{joypad_dpad} and
\fname{joypad_buttons} are ever used; the upper nibble is not part of
the interface.  Section~\ref{controls} has the correct table.

\subsection{Serial output}
\label{serial}

There is no serial shift register.  Two intercept paths turn a game's
output into text, and both are what makes automated testing possible.

\begin{description}

\item[Register path --- \texttt{0xFF02}.]  Writing the exact byte
  \cs{0x81} (internal clock plus start) prints the byte previously
  written to \fname{SB} at \texttt{0xFF01} with \cs{putchar} and
  \cs{fflush}, and then stores \cs{0x01} in \fname{SC} so that a ROM
  polling bit~7 of \cs{SC} terminates.  This is the protocol used by
  mooneye and SameSuite: the ROM prints a pass signature of
  \cs{03 05 08 0D 15 22} or six \cs{0x42} bytes.

\item[RAM path --- \texttt{0xA004} and above.]  Writes to cartridge RAM
  at or above \texttt{0xA004} print every printable byte
  (\cs{0x20}--\cs{0x7E}, newline, carriage return, tab) to standard
  output, and append \emph{every} byte --- including NUL terminators ---
  to \fname{/tmp/blargg_raw.bin}.  This is Bla'rg's
  \fname{shell.s} text-output protocol.

\end{description}

\marginlabel{Raw trace:}
\fname{/tmp/blargg_raw.bin} is opened once per process and truncated at
first use.  It is the way to recover a diagnostic that contains
non-printable bytes, and it is compiled out of the ESP32 build.

The RAM path is only reached by carts whose mapper does not return early
--- MBC1, MBC5 and ROM-only.  MBC2 and MBC3 returns skip it, so text
output from those mappers is lost.

\subsection{OAM DMA}
\label{oam-dma}

Writing a source page to \texttt{0xFF46} starts a transfer of 160 bytes
into the sprite attribute table.  The timing is explicit:

\begin{longtable}{@{}C{0.22\textwidth}L{0.68\textwidth}@{}}
\toprule
\textbf{When} & \textbf{What happens}\\
\midrule
\endhead
Write &
  \cs{dma\_active = true}, \cs{dma\_start\_delay = 2},
  \cs{dma\_offset = 0}, source page latched\\
+1 M-cycle & Delay decrements; OAM is still accessible\\
+2 M-cycle & Delay reaches zero, the first byte copies --- OAM is now
  blocked\\
+3 to +161 M-cycles & One byte per M-cycle, 159 more bytes\\
+162 M-cycle & \cs{dma\_offset == 0xA0}; the transfer ends and OAM
  reopens\\
\bottomrule
\end{longtable}

While a transfer is copying, reads from \texttt{0xFE00--0xFE9F} return
\cs{0xFF} and writes are dropped.  A write to \texttt{0xFF46} issued while
bytes are already moving is queued: the current transfer continues and a
restart begins two machine cycles later.  A write during the initial
two-cycle delay, or while a restart is already pending, restarts
immediately.

The source address is reconstructed per byte as
\cs{(page << 8) | (offset \& 0xFF)} and read through the ordinary bus.
It is neither masked nor validated, so a source page of \cs{0xFE} copies
\fname{0xFF} bytes --- because OAM reads are themselves blocked during
the copy --- and pages outside the \cs{0xXX00--0xXX9F} window are
permitted.  Section~\ref{dma-limits} lists this among the known gaps.

\subsection{\cs{IF}, \cs{IE} and the interrupt flags}

Covered in section~\ref{interrupts}: \fname{IF} lives in
\cs{bus->io[0x0F]}, \cs{IE} in \cs{bus->ie}, and hardware request flags
live in the components that raise them.  All three are included in save
states; the transient \cs{joypad\_interrupt} flag is not.

\subsection{The CGB speed switch}
\label{key1}

\fname{KEY1} at \texttt{0xFF4D} is the only CGB register the emulator
implements:

\begin{itemize}
\item Read: \cs{(double\_speed ? 0x80 : 0x00) | (io[0x4D] \& 0x01)} ---
      bit~7 reports the \emph{current} mode, bit~0 the prepared one.
\item Write: only bit~0 is stored.
\item \cs{STOP} performs the actual switch, flipping
      \cs{bus->double\_speed} and clearing the request bit.
\end{itemize}

Once \fname{double_speed} is set, \fname{gb_machine_tick} divides the
cycle count handed to the PPU and the APU by two, and \fname{gb_loop}
counts dots instead of T-cycles.  The frame budget is therefore identical
in both modes, as on hardware.  The rest of the CGB --- palettes, VRAM
banking, \cs{GDMA}/\cs{HDMA}, the boot ROM --- is listed in
section~\ref{cgb-limits}.

\clearpage

\clearpage

% ===========================================================================
\section{Memory bank controllers}
\label{mappers}

\fname{src/bus.c} implements four mapper families, selected purely by the
cartridge type byte at header offset \cs{0x0147}.  There is no fall-back
mapper: an unknown type is treated as ROM-only, and the type is stored as
zero so that the persistence layer can tell the difference.

\begin{longtable}{@{}L{0.14\textwidth}L{0.30\textwidth}L{0.44\textwidth}@{}}
\toprule
\textbf{Types} & \textbf{Mapper} & \textbf{Identification predicate}\\
\midrule
\endhead
\cs{0x01--0x03} & MBC1 & \cs{type >= 0x01 \&\& type <= 0x03}\\
\cs{0x05--0x06} & MBC2 & \cs{type >= 0x05 \&\& type <= 0x06}\\
\cs{0x0F--0x13} & MBC3 & \cs{type >= 0x0F \&\& type <= 0x13}\\
\cs{0x19--0x1E} & MBC5 & \cs{type >= 0x19 \&\& type <= 0x1E}\\
otherwise & none & \cs{mbc\_type} is stored as 0\\
\bottomrule
\end{longtable}

Ranges rather than equality tests are used deliberately, so that the
rumble and battery variants (\cs{0x1B}--\cs{0x1E}) are covered without a
separate case each.  Every predicate is inlined into the read and write
paths, which keeps \fname{bus_read_rom} and \fname{bus_read_sram} free
of dispatch tables.

\marginlabel{Addressing:}
All four mappers see the same 32\,KiB window.  A ROM read is always
\fname{rom[bank * 0x4000 + (addr \& 0x3FFF)]}, and every mapper's
resolution of \cs{bank} is finally clamped with
\cs{bank \& (rom\_banks - 1)}.  Because \cs{rom\_banks} is always a power
of two chosen by \fname{gb_bus_load_rom}, that clamp can never
out-of-bounds, whatever the cartridge claims in its header.

\subsection{MBC1}

MBC1 has two mode-dependent behaviours and one special case, all confined
to \fname{mbc1_bank_for_addr}.

\begin{longtable}{@{}C{0.24\textwidth}C{0.34\textwidth}C{0.36\textwidth}@{}}
\toprule
\textbf{Window} & \textbf{ROM mode} & \textbf{RAM mode}\\
\midrule
\endhead
\cs{0x0000--0x3FFF} & always bank~0 & bank \cs{<< 5} from the 2-bit
  register\\
\cs{0x4000--0x7FFF} & \cs{rom\_bank | (ram\_bank << 5)} &
  \emph{identical} to ROM mode\\
\midrule
Cartridge RAM & bank~0 & the 2-bit register\\
\bottomrule
\end{longtable}

Three hardware details that the code reproduces:

\begin{description}

\item[Bank-0 substitution] \cs{0x00}, \cs{0x20}, \cs{0x40} and \cs{0x60}
      are replaced by \cs{0x01}, \cs{0x21}, \cs{0x41} and \cs{0x61} in the
      switchable window, so bank~0 is only reachable through the fixed
      region.
\item[The mode bit only affects the fixed region] The upper two ROM bank
      bits are OR-ed in in \emph{both} modes; the mode select at
      \cs{0x6000} changes only what \cs{0x0000--0x3FFF} shows and which
      cartridge RAM bank is addressed.  This matches real MBC1B1
      hardware, which is not what the classic Game Boy manuals describe.
\item[MBC1M multicarts] A 1\,MiB MBC1 cart may actually be a multicart with
      four 256\,KiB games.  The header cannot tell you, so
      \fname{gb_bus_load_rom} computes CRC-32 over the 48-byte
      Nintendo logo at offset \cs{0x0104} in each of the four pages and
      sets \cs{mbc1\_multicart} when at least three match the reference
      value \cs{0x46195417} --- the same heuristic mooneye-gb uses.  In
      that mode the ROM bank register contributes only its low four bits
      and the 2-bit register selects the game, with bank~0 valid.
\end{description}

The four registers occupy the standard windows:

\begin{longtable}{@{}L{0.26\textwidth}C{0.22\textwidth}L{0.42\textwidth}@{}}
\toprule
\textbf{Window} & \textbf{Field} & \textbf{Effect}\\
\midrule
\endhead
\texttt{0x0000--0x1FFF} & RAM enable & Set when the low nibble of the
  written byte is \cs{0x0A}\\
\texttt{0x2000--0x3FFF} & ROM bank & Low five bits
  (\cs{mbc1\_rom\_bank})\\
\texttt{0x4000--0x5FFF} & 2-bit register & Upper ROM bank bits \emph{or}
  the RAM bank, depending on the mode\\
\texttt{0x6000--0x7FFF} & Mode & \cs{0} = ROM banking,
  \cs{1} = RAM banking\\
\bottomrule
\end{longtable}

Cartridge RAM reads return \cs{0xFF} while RAM is disabled; the read path
checks \cs{mbc1\_ram\_enable} before anything else.

\subsection{MBC2}

MBC2 is the odd one out: it has 512 half-bytes of RAM on the cartridge
itself, addressable through the same \cs{0xA000--0xBFFF} window with the
low nine bits of the address, and only four bits per byte are
meaningful.

\begin{itemize}
\item Writes to \cs{0xA000--0xBFFF} store \cs{value \& 0x0F} and return
      immediately --- which means the serial-output interception of
      section~\ref{serial} never runs for MBC2 games.
\item Reads return \cs{0xF0 | (ram[addr \& 0x1FF] \& 0x0F)} when RAM is
      enabled and \cs{0xFF} when it is not.  The upper nibble reads as
      \cs{0xF} because the register only has four RAM bits; this is what
      the mooneye \fname{bits_unused} test checks.
\item The mapper registers live in the \emph{fixed} window and are
      selected by address bit~8: a write with bit~8 clear is the RAM
      enable (\cs{0x0A} in the low nibble), a write with bit~8 set is the
      4-bit ROM bank.
\item Writes to \cs{0x4000--0x7FFF} are ignored outright, which the
      mooneye test suite verifies on MBC2A hardware.
\item Bank~0 is substituted by bank~1 in the switchable window, and the
      result is masked to the cartridge size.
\end{itemize}

\attention
MBC2 is the one mapper whose 512-byte save file is a different size from
the 8\,KiB-or-more allocation the loader makes for every other cart
(section~\ref{sram-sizing}), and \fname{gb_battery_save} special-cases
it accordingly.

\subsection{MBC3}

MBC3 adds a real-time clock and a register-select trick: the same 2-bit
window that selects a RAM bank can instead select one of five RTC
registers.

\begin{longtable}{@{}L{0.26\textwidth}C{0.20\textwidth}L{0.44\textwidth}@{}}
\toprule
\textbf{Window} & \textbf{Field} & \textbf{Effect}\\
\midrule
\endhead
\texttt{0x0000--0x1FFF} & RAM/RTC enable & Low nibble \cs{0x0A} enables
  both; anything else disables both\\
\texttt{0x2000--0x3FFF} & ROM bank & Seven bits; a written 0 becomes 1
  immediately, so bank~0 is unreachable here\\
\texttt{0x4000--0x5FFF} & RAM bank / RTC select & \cs{0}--\cs{3} select a
  RAM bank, \cs{0x08}--\cs{0x0C} select the RTC registers\\
\texttt{0x6000--0x7FFF} & Latch & Writing \cs{0x00} then \cs{0x01}
  latches the clock\\
\bottomrule
\end{longtable}

The RTC registers are kept in \cs{bus->mbc3\_seconds} (0--59),
\cs{minutes} (0--59), \cs{hours} (0--23), a 9-bit
\cs{day\_counter}, a \cs{day\_carry} flag and a \cs{timer\_halt} flag.
Writes are range-checked by modulo (\cs{value \% 60},
\cs{value \% 24}), which is what the RTC register behaviour amounts to
for out-of-range values.

Reads and writes to RTC registers are reached through the normal RAM
window after selecting them, and they set \cs{sram\_dirty} so that they
end up in the battery save (section~\ref{battery-format}).

\attention
The latch sequence only advances a state machine; it does not copy the
host clock into the registers.  There is a \cs{TODO} marker for that.  A
game that latches before reading will read the counters it found at load
time, which is stable but not wall-clock accurate.

\subsection{MBC5}

MBC5 is the simplest of the four --- nine ROM bank bits, four RAM bank
bits, no mode bit, no bank-0 substitution --- and it is the only mapper
that can address a 2\,MiB cartridge.

\begin{longtable}{@{}L{0.26\textwidth}C{0.22\textwidth}L{0.42\textwidth}@{}}
\toprule
\textbf{Window} & \textbf{Field} & \textbf{Effect}\\
\midrule
\endhead
\texttt{0x0000--0x1FFF} & RAM enable & Low nibble \cs{0x0A}\\
\texttt{0x2000--0x2FFF} & ROM bank low & All eight bits, no masking\\
\texttt{0x3000--0x3FFF} & ROM bank high & Bit~0 only, giving bank
  numbers up to 511\\
\texttt{0x4000--0x5FFF} & RAM bank & Low four bits\\
\texttt{0x6000--0x7FFF} & Rumble & Stored nowhere; ignored\\
\bottomrule
\end{longtable}

\attention
Because MBC5 does not substitute bank~1 for a bank~0 write, the loader
starts an MBC5 cartridge with bank~1 already selected
(\cs{mbc5\_rom\_bank\_l = 1}).  Without this, the mooneye harness --- which
copies its \fname{memcpy} and friends out of bank~1 into WRAM before
switching banks --- would fault on its first library call.

\subsection{Cartridge RAM sizing}
\label{sram-sizing}

The allocation comes from header byte \cs{0x0149}, with a floor:

\begin{longtable}{@{}C{0.18\textwidth}C{0.26\textwidth}C{0.46\textwidth}@{}}
\toprule
\textbf{Header} & \textbf{Allocated} & \textbf{Remark}\\
\midrule
\endhead
\cs{0x00} & 8\,KiB & No RAM declared, but still allocated\\
\cs{0x01} & 8\,KiB & 2\,KiB requested, raised to the floor\\
\cs{0x02} & 8\,KiB & The only size real carts use here\\
\cs{0x03} & 32\,KiB &\\
\cs{0x04} & 128\,KiB &\\
other & 8\,KiB & Treated as 8\,KiB\\
\bottomrule
\end{longtable}

\attention
The 8\,KiB floor exists for the test suites: Bla'rg's \fname{shell.s}
buffers its serial output in SRAM at \cs{0xA004} and above, so a cartridge
that declares no RAM still needs a writable window there or the diagnostic
never reaches standard output.  The trade-off is that
\cs{sram\_banks} is 1 for most carts, so the bank mask is a no-op.

\seealso
\cs{rom\_banks} is computed the same way --- the next power of two at or
above the number of 16\,KiB banks the file actually occupies, minimum 2
(section~\ref{rom-loading}).

\clearpage

% ===========================================================================
\section{The timer and divider}
\label{timer}

The timer is the one component that does \emph{not} get its cycle count
scaled.  \fname{gb_machine_tick} hands the raw T-cycle count to
\fname{gb_timer_step} and the scaled dot count to the PPU and the APU,
because \cs{DIV} is derived from the CPU clock, not from the dot clock.
The consequence is that in CGB double-speed mode \cs{DIV} increments twice
as fast in wall-clock terms, exactly as on hardware.

\subsection{The divider}

\fname{gb_timer.internal_counter} is a 16-bit free-running counter
incremented once per T-cycle.  \fname{DIV} is simply its upper byte, and
\fname{gb_timer_init} seeds it with \cs{0xABCC} so that the first
\fname{DIV} read returns the post-boot \cs{0xAB}.

The timer input is a \emph{bit} of that counter selected by \cs{TAC}, and
\fname{TIMA} increments on the falling edge of that bit --- the standard
hardware model, implemented by sampling the bit before and after the
increment:

\begin{Code}
bool bit_before = get_timer_bit(timer);   // TAC enable + selected bit
timer->internal_counter++;
bool bit_after  = get_timer_bit(timer);
if (bit_before && !bit_after) {
    if (tima == 0xFF) { tima = 0; reload_pending = true; request_interrupt(); }
    else              tima++;
}
\end{Code}

\marginlabel{Four rates:}
\fname{TAC} bits 1--0 select the counter bit, and therefore the
\fname{TIMA} increment period in T-cycles:

\begin{longtable}{@{}C{0.14\textwidth}C{0.20\textwidth}C{0.20\textwidth}C{0.34\textwidth}@{}}
\toprule
\textbf{\cs{TAC} bits} & \textbf{Counter bit} & \textbf{Period} &
\textbf{Frequency}\\
\midrule
\endhead
\cs{00} & 9 & 1024 T-cycles & 4096\,Hz\\
\cs{01} & 3 & 16 T-cycles & 262\,144\,Hz\\
\cs{10} & 5 & 64 T-cycles & 65\,536\,Hz\\
\cs{11} & 7 & 256 T-cycles & 16\,384\,Hz\\
\bottomrule
\end{longtable}

Bit~2 of \cs{TAC} is the enable.  \fname{gb_get_timer_bit} returns
\cs{false} immediately when it is clear, which means writing a zero to
\fname{TAC} can itself produce a falling edge --- and therefore increment
\fname{TIMA}.  That is real hardware behaviour, and the code lets it
happen rather than special-casing the register write.

\subsection{The \cs{TIMA} overflow delay}

An overflow does not reload \cs{TIMA} immediately.  \fname{tima} is set to
\cs{0x00} and \cs{tima\_reload\_pending} is raised; the reload happens on
the \emph{next} call to \fname{gb_timer_step}, before the counter is
incremented:

\begin{Code}
if (timer->tima_reload_pending) {
    timer->tima = timer->tma;
    timer->tima_reload_pending = false;
}
\end{Code}

That one-cycle delay is what makes the mooneye
\fname{timer/tima_reload} test pass, and it gives a game exactly one
cycle in which \fname{TIMA} reads as zero.  A write to \fname{TIMA} during
that window cancels the pending reload, again matching hardware.

The interrupt request is raised at the moment of the overflow, not at the
reload, so the flag can already be in \cs{IF} one T-cycle before
\fname{TIMA} reads \cs{TMA}.

\subsection{Register behaviour}

\begin{longtable}{@{}L{0.20\textwidth}L{0.74\textwidth}@{}}
\toprule
\textbf{Register} & \textbf{Behaviour}\\
\midrule
\endhead
\fname{DIV} (\cs{0xFF04}) &
  Read: upper byte of the internal counter.  Write: \emph{any} value
  clears the whole counter, and --- if the timer was enabled and the
  selected bit was 1 --- that reset counts as a falling edge and
  increments \fname{TIMA}, overflow included.  Both behaviours are
  required by the mooneye \fname{div_write} and \fname{div_trigger}
  tests.\\
\fname{TIMA} (\cs{0xFF05}) &
  Read and write directly.  Writing cancels a pending reload.\\
\fname{TMA} (\cs{0xFF06}) &
  Read and write directly; the value is latched into \fname{TIMA} when the
  pending reload fires.\\
\fname{TAC} (\cs{0xFF07}) &
  Only bits 0--2 are writable.  Reads return \cs{tac | 0xF8}, so the four
  unused bits read as one.\\
\bottomrule
\end{longtable}

\seealso
\fname{internal_counter} starts at \cs{0xABCC} and advances one step per
T-cycle, so \fname{gb_timer_step(timer, n)} is the \emph{only} thing that
drives it.  Calling the PPU or the APU directly, without the timer,
produces a machine whose \cs{DIV} does not move.

\clearpage

% ===========================================================================
\section{The PPU}
\label{ppu}

The PPU is a scanline renderer, not a pixel-by-pixel one.  It advances a
dot counter, changes modes at fixed thresholds, and draws a whole
scanline into the frame buffer at a single moment --- the mode~3 to
mode~0 transition.  There is no FIFO, no mid-scanline fetch and no
per-pixel timing.

\subsection{Registers}

\begin{longtable}{@{}C{0.14\textwidth}L{0.26\textwidth}L{0.50\textwidth}@{}}
\toprule
\textbf{Address} & \textbf{Name} & \textbf{Bits and behaviour}\\
\midrule
\endhead
\cs{0xFF40} & \cs{LCDC} &
  Bit~7 LCD enable, 6 window map, 5 window enable, 4 tile data source,
  3 tile map source, 2 sprite size, 1 sprite enable, 0 BG/window
  enable.\\
\cs{0xFF41} & \cs{STAT} &
  Bits 1--0 are the mode (hardware-owned), bit~2 the \cs{LYC=LY}
  coincidence flag, bits~3--5 the mode-interrupt enables, bit~6 the
  \cs{LYC=LY} interrupt enable.  Writes keep only bits 3--6; reads
  additionally return \cs{0x80}.\\
\cs{0xFF42} & \cs{SCY} & Background scroll Y\\
\cs{0xFF43} & \cs{SCX} & Background scroll X\\
\cs{0xFF44} & \cs{LY} & Current scanline; read-only, writes ignored\\
\cs{0xFF45} & \cs{LYC} & The scanline compared against \cs{LY}\\
\cs{0xFF47} & \cs{BGP} & Two bits per shade for the background\\
\cs{0xFF48} & \cs{OBP0} & As above for sprites using palette 0\\
\cs{0xFF49} & \cs{OBP1} & As above for sprites using palette 1\\
\cs{0xFF4A} & \cs{WY} & First scanline on which the window is visible\\
\cs{0xFF4B} & \cs{WX} & Window left edge, in the hardware's
  \cs{X+7} encoding\\
\bottomrule
\end{longtable}

Registers outside \cs{0xFF40--0xFF4B} read as \cs{0xFF}.  In particular the
CGB palette registers \cs{0xFF68} and \cs{0xFF69} are not decoded and fall
into the bus's shadow array, so they read back as whatever was last
written there --- zero, initially.

\subsection{Mode state machine}
\label{ppu-modes}

\fname{gb_ppu_step} receives a dot count, adds it to \cs{ppu->dots}, and
checks the mode stored in the low two bits of \fname{STAT}:

\begin{longtable}{@{}C{0.14\textwidth}L{0.24\textwidth}C{0.20\textwidth}L{0.32\textwidth}@{}}
\toprule
\textbf{Mode} & \textbf{Name} & \textbf{Dots} & \textbf{Transition at}\\
\midrule
\endhead
\cs{2} & OAM search & 160 & \cs{dots \>= 160} $\rightarrow$ mode~3\\
\cs{3} & Transfer & 344 more & \cs{dots \>= 504} $\rightarrow$ draw, mode~0\\
\cs{0} & H-blank & 408 more & \cs{dots \>= 912} $\rightarrow$ next line\\
\cs{1} & V-blank & 912 per line & \cs{ly} 144--153, then wrap to 0\\
\bottomrule
\end{longtable}

One scanline is therefore 912 dots, and a whole frame is
\cs{144 * 912} dots of active lines plus ten blanking lines.  The drawing
itself happens inside the mode~3 to mode~0 transition, in this order:
background, window, sprites.

Two events are raised by the state machine:

\begin{itemize}
\item Entering mode~1 sets \cs{frame\_ready} and \cs{vblank\_interrupt}.
      \cs{frame\_ready} is what tells \fname{gb_loop} to call
      \cs{render}.
\item Every mode transition consults the mode-interrupt enables, and
      every line boundary consults \fname{check_lyc_stat_interrupt},
      which maintains the coincidence flag and raises
      \cs{stat\_interrupt} when bit~6 is set.
\end{itemize}

\attention
Both interrupt flags are \emph{booleans, not bits in \cs{STAT}}.  They are
latched into \cs{IF} by \fname{gb_handle_interrupts} at the next
instruction boundary and cleared there.  A game that clears \cs{STAT}'s
enable bits after the mode has already been entered will not see the
interrupt, which differs from hardware but only matters for code that
races the transition.

\subsection{Turning the LCD on and off}

Writing \cs{LCDC} bit~7 from 1 to 0 or from 0 to 1 resets the PPU:

\begin{itemize}
\item \emph{Off:} \cs{ly = 0}, \cs{dots = 0}, mode forced to 0
      (H-blank), \cs{first\_line\_after\_enable} cleared.  \fname{gb_ppu_step}
      returns immediately from then on, so the display is frozen on the
      last rendered frame.
\item \emph{On:} the same three writes, plus
      \cs{first\_line\_after\_enable = true}, which produces a shortened
      first line of 900 dots.  At \cs{dots \>= 900} the PPU jumps
      straight to line~1 in mode~2.
\end{itemize}

The 900-dot first line exists because the bl'rg \fname{oam_bug}
composite reads \cs{LY} at roughly dot 896 and again at dot 904 and
expects to see 0 and then~1.  Getting that wrong is one of the ways the
suite fails.

\subsection{Rendering a scanline}

\fname{ppu_render_bg} walks all 160 screen columns and, for each one,
computes the source address and writes one pixel:

\begin{Code}
line_y    = ly + SCY;
line_x    = x  + SCX;
tile_row  = (line_y / 8) * 32;
tile_col  = line_x / 8;
map_addr  = tile_map_base + tile_row + tile_col;
tile      = bus_read(map_addr);
tile_addr = tile_data_base + tile*16 + (line_y % 8)*2;
bit       = 7 - (line_x % 8);
color     = ((plane2 >> bit & 1) << 1) | (plane1 >> bit & 1);
buffer[ly][x] = palette(BGP, color);
\end{Code}

Both addressing modes of \cs{LCDC} bit~4 are implemented: unsigned from
\cs{0x8000} and signed from \cs{0x9000} with the tile index read as a
\fname{int8_t}.  Every fetch goes through \fname{gb_bus_read}, so the
renderer sees the same memory the CPU does --- including mapper effects on
a game that (wrongly) puts tile data in ROM.

\fname{ppu_render_window} is the same loop with \cs{line\_y =
ly - WY}, \cs{line\_x = x - (WX - 7)}, the LCDC bit~6 map, and a
skip of every column left of the window's edge.  The window has no
internal line counter: its position within the map is recomputed from
\cs{LY - WY} every line, which is why \cs{WY} must not change mid-frame
if you want hardware-identical results.

\fname{ppu_render_sprites} collects at most ten sprites whose 8 or 16
pixel band contains the current line, then draws them in reverse
collection order so that the lowest OAM index ends up on top:

\begin{longtable}{@{}L{0.24\textwidth}L{0.68\textwidth}@{}}
\toprule
\textbf{OAM byte} & \textbf{Meaning}\\
\midrule
\endhead
0 & Screen Y, biased by \cs{-16}; a sprite spans \cs{ly} in
  \cs{[y-16, y-16+height)}\\
1 & Screen X, biased by \cs{-8}\\
2 & Tile index; bit~0 is ignored in 8$\times$16 mode\\
3 & Bit~7 BG priority, bit~6 Y flip, bit~5 X flip, bit~4 palette\\
\bottomrule
\end{longtable}

Colour index 0 is transparent, so sprites never erase the background.
Only the first ten matching sprites are drawn, which is the hardware's
per-line limit rather than a choice.

\attention
The BG-priority bit (bit~7 of OAM byte~3) is parsed but not applied:
\fname{ppu_render_sprites} reads it into \cs{bg\_over} and then, as the
in-code comment says, cannot tell whether the pixel underneath is
background colour 0 without keeping the background colour indices around.
Sprites therefore always draw over the background, and over each other, in
OAM order.  This is visible in games that put a sprite behind a background
tile and expect it to disappear.

\subsection{Palettes}
\label{ppu-palettes}

There is one palette conversion routine for the whole system,
\fname{ppu_get_color(palette_reg, color_index)}, and one fixed
four-entry table:

\begin{longtable}{@{}C{0.24\textwidth}C{0.34\textwidth}L{0.30\textwidth}@{}}
\toprule
\textbf{Shade} & \textbf{Colour} & \textbf{Value}\\
\midrule
\endhead
0 & white & \cs{0xFFFFFFFF}\\
1 & light grey & \cs{0xFFAAAAAA}\\
2 & dark grey & \cs{0xFF555555}\\
3 & black & \cs{0xFF000000}\\
\bottomrule
\end{longtable}

The shade is \cs{(palette\_reg >> (color\_index * 2)) \& 3}, so the DMG
palette registers are honoured exactly; only the CGB's per-tile palettes
are missing.  Because the frame buffer is ARGB8888 with an opaque alpha,
the terminal and WASM front ends can read the pixel as a plain
\fname{uint32_t} and shift it, and the ESP32 front end can convert it to
RGB565 without knowing anything about the Game Boy.

\subsection{What the PPU does not do}
\label{ppu-limits}

The following are absent, and are the reason several bl'rg and SameSuite
tests fail:

\begin{itemize}
\item No per-pixel fetch timing, no FIFO, and no per-cycle object
      priority resolution.
\item No blocking of VRAM access during modes 2 and 3, or of OAM during
      modes 2 and 3, or of writes to \cs{LY}.
\item No \cs{STAT} blocking flags (the bit~7 ``STAT blocking'' bit), so a
      write to \cs{STAT} that hardware would delay takes effect at once.
\item No DMG OAM scan bug.  Sprites are read straight out of OAM at draw
      time, whereas hardware latches them during mode~2 and can be
      corrupted by writes in the meantime.
\item No CGB: no palette RAM, no tile attribute byte, no VRAM banking, no
      separate window line counter, no \cs{BGPI}/\cs{BGPD}.
\end{itemize}

See section~\ref{known-gaps} for the full list and
section~\ref{testing} for how each gap shows up in the suites.

\clearpage

% ===========================================================================
\section{The APU}
\label{apu}

The APU synthesises all four channels at the emulated sample rate and
hands the result to the host through a lock-free ring buffer.  This is
one of the more complete parts of the emulator; the gaps that remain are
listed at the end of the section.

\subsection{Registers}

All forty-eight sound registers live in one array, \fname{apu->regs[0x30]},
indexed by \cs{addr - 0xFF10}, which is also how wave RAM is stored: the
sixteen bytes at \cs{0xFF30}--\cs{0xFF3F} occupy
\cs{regs[0x20]}--\cs{regs[0x2F]}.

\begin{longtable}{@{}L{0.14\textwidth}L{0.20\textwidth}L{0.56\textwidth}@{}}
\toprule
\textbf{Offset} & \textbf{Register} & \textbf{Role}\\
\midrule
\endhead
\cs{0x00} & \cs{NR10} & Sweep period, negate, shift for channel~1\\
\cs{0x01} & \cs{NR11} & Duty and length load for channel~1\\
\cs{0x02} & \cs{NR12} & Envelope and DAC enable for channel~1\\
\cs{0x03} & \cs{NR13} & Frequency low byte, channel~1\\
\cs{0x04} & \cs{NR14} & Frequency high bits, length enable, trigger\\
\cs{0x05}--\cs{0x09} & \cs{NR21}--\cs{NR24} & As above for channel~2\\
\cs{0x0A}--\cs{0x0E} & \cs{NR30}--\cs{NR34} & DAC enable, sample volume,
  frequency, length enable, trigger for channel~3\\
\cs{0x0F} & \cs{NR41} & Length load for channel~4\\
\cs{0x10} & \cs{NR42} & Envelope and DAC enable for channel~4\\
\cs{0x11} & \cs{NR43} & Shift clock, width mode, divisor\\
\cs{0x12} & \cs{NR44} & Length enable and trigger\\
\cs{0x14} & \cs{NR50} & Master volume, left (bits 6--4) and right (2--0)\\
\cs{0x15} & \cs{NR51} & Panning: bits 0--3 right, bits 4--7 left\\
\cs{0x16} & \cs{NR52} & Master power (bit~7) and the four channel-on flags
  (bits 0--3)\\
\cs{0x20}--\cs{0x2F} & wave RAM & 32 four-bit samples\\
\bottomrule
\end{longtable}

Reads apply a per-register mask so that unused bits read as one, exactly
as on hardware; \fname{NR52} is special-cased to synthesise the channel
flags from \cs{ch\_on[]}, and wave RAM is returned unmasked.

\subsection{Power}

\fname{NR52} bit~7 is \cs{apu->power}.

\begin{description}

\item[Powering off] Clears \cs{NR10}--\cs{NR51}, silences and disables
      every channel, disables every length counter, zeroes the frequency
      timers, the duty positions, the wave position, the frame sequencer
      and the LFSR, and reloads the LFSR with \cs{0x7FFF}.  On a CGB cart
      the length counters are zeroed as well; on a DMG they are left
      alone.
\item[Powering on] Sets \cs{power} and sets \cs{frame\_div} to 1024, which
      places the sequencer in the second half of its period --- the
      documented behaviour for a power-on during a frame.
\item[While powered off] \cs{NR10}--\cs{NR51} are read-only.  On a DMG
      cart the one exception is that a write to \cs{NRx1} still updates the
      length counter, which is what the length tests check.  Wave RAM
      remains writable in both cases.
\end{description}

\subsection{The frame sequencer}

\fname{apu->frame_div} accumulates dots and steps every 2048 of them,
i.e.\ at 512\,Hz in normal speed and 1024\,Hz in double speed.  The eight
steps do the following:

\begin{longtable}{@{}C{0.14\textwidth}L{0.76\textwidth}@{}}
\toprule
\textbf{Step} & \textbf{Action}\\
\midrule
\endhead
0, 2, 4, 6 & Clock the length counters of all four channels; a counter
  reaching zero switches its channel off\\
2, 6 & Step the sweep unit of channel~1\\
7 & Step the volume envelopes of channels 1, 2 and 4\\
\bottomrule
\end{longtable}

\marginlabel{Why 2048 dots:}
The APU is stepped in dots, and the sequencer is defined in terms of the
CPU frequency.  8192 dots --- 2048 T-cycles at normal speed --- is one
512\,Hz period, so the divisor is correct in both speed modes without any
scaling.

\subsection{Frequency timers}

Each channel has a countdown that is decremented by the dot count every
step; while it is non-positive the channel advances and the timer is
reloaded.  The reload value is \cs{(2048 - frequency) * 4} for the two
square channels and the wave channel, where \cs{frequency} is the
eleven-bit value assembled from \cs{NRx3} and the low three bits of
\cs{NRx4}.  For the noise channel it is four times one of eight divisors
(\cs{8, 16, 32, 48, 64, 80, 96, 112}) taken from the low three bits of
\cs{NR43}.

What advances on each expiry:

\begin{itemize}
\item \emph{Channel 1 and 2:} \cs{duty\_pos} steps modulo 8.
\item \emph{Channel 3:} \cs{wave\_pos} steps modulo 32.
\item \emph{Channel 4:} the 15-bit LFSR shifts right, feeding bit 0 into
      bit 14.  With \cs{NR43} bit~3 set, bit~6 receives the XOR feedback
      instead of bit~1, which is the 7-bit width mode.
\end{itemize}

Duty patterns are the four standard constants --- \cs{0x01},
\cs{0x81}, \cs{0x87}, \cs{0x7E} --- indexed by bits 6--7 of \cs{NRx1}.
The output of a square channel is its current volume when the selected
bit of the duty pattern is set and zero otherwise.

\subsection{Volume envelopes}

\fname{vol[3]}, \cs{env\_period[3]}, \cs{env\_counter[3]} and
\cs{env\_dir[3]} cover channels 1, 2 and 4; the wave channel has no
envelope.  At step~7 each channel with a non-zero period decrements its
counter and, when that counter expires, reloads it and moves the volume
one step up (direction set) or down (direction clear), clamped to 0--15.
A triggering event reloads the volume from the high nibble of \cs{NRx2}
and resets the counter.

\attention
The envelope parameters are read out of \cs{NRx2} only on trigger, and
\cs{env\_period} and \cs{env\_dir} are never updated from a later write to
\cs{NRx2}.  Hardware applies them immediately.  In practice a game writes
\fname{NRx2} once and then triggers, so this is only visible to tests
that rewrite the envelope mid-note.

\subsection{Sweep}

Channel~1's frequency is shadowed in \cs{sweep\_freq}.  At sequencer
steps 2 and 6 the counter is decremented, and when it expires the
frequency is recomputed as \cs{sweep\_freq} plus or minus
\cs{(sweep\_freq >> shift)}, where the sign is \cs{sweep\_negate}.  If the
result exceeds 2047
the channel is switched off and the sweep disabled; otherwise the new
value is written back into \cs{NR13} and \cs{NR14} and the shadow is
recalculated.  A second overflow check runs after the rewrite.

A trigger reloads the shadow from \cs{NR13}/\cs{NR14}, sets the counter
from \cs{NR10} bits 4--6, and --- if the shift is non-zero --- performs
the overflow calculation immediately, which is how the ``channel off on
trigger when the sweep overflows'' case is handled.  Writing \cs{NR10}
with the negate bit cleared after negation has been used switches the
channel off, via \cs{sweep\_negate\_used}.

\subsection{Length counters}

Writes to the four \cs{NRx1} registers load the counters:
\cs{64 - (value \& 0x3F)} for the square and noise channels and
\cs{256 - value} for the wave channel.  Bit~6 of each \cs{NRx4} is the
length enable.  On trigger, a counter that is currently zero is
unfrozen at 64 (or 256), and if the length enable is set and the
sequencer is in the first half of its period, an extra length clock is
applied immediately.  The same extra clock is applied when the length
enable bit goes from 0 to 1 during that half-period.

\subsection{Triggers and DAC}

A trigger (bit~7 of \cs{NRx4}) switches the channel on, reloads its
frequency timer, resets the duty position or the wave position or the
LFSR, and reloads the envelope.  Immediately afterwards the DAC is
checked: channels~1, 2 and 4 need \cs{(NRx2 \& 0xF8) != 0} and channel~3
needs \cs{NR30} bit~7, so a trigger with the DAC disabled silently leaves
the channel off.  Conversely, writing a zero upper five bits to \cs{NRx2}
switches the channel off without a trigger.

\subsection{Mixing and output}

Every \cs{GB\_CPU\_FREQ / GB\_APU\_SAMPLE\_RATE} dots --- 190 at 44.1\,kHz
--- the four channels are sampled:

\begin{enumerate}
\item Each channel is normalised to \cs{[0, 1]}.  The wave channel reads
      the sample nibble selected by \cs{wave\_pos}, where the high nibble
      of a wave RAM byte comes first, and shifts it right by
      \cs{NR30} bits 5--4 (0 mutes, 1 is full scale, 2 halves it, 3
      quarters it).
\item \cs{NR51} routes each channel to the left and/or the right bus.
\item \cs{NR50} scales each side by its three-bit volume divided by~7.
\item The two sides are added and divided by four --- the code divides by
      8, which leaves the result at roughly half scale --- clamped to
      \cs{[-1, 1]} and pushed onto the ring buffer.
\end{enumerate}

The buffer is a 4096-entry single-producer single-consumer ring of
\fname{float}s with \cs{buf\_write} and \cs{buf\_read} marked
\cs{volatile}.  The producer is \fname{gb_apu_step} on the emulation
thread; the consumer is \fname{gb_apu_buf_pop}, which the SDL audio
callback and the ESP32 audio task call.  If the consumer falls behind, the
producer simply drops samples --- the push is skipped when the next index
would collide with the read index, rather than overwriting unread audio.

\subsection{Power-on defaults}

\fname{gb_apu_init} clears everything and then sets three registers that
matter: \cs{NR52} to \cs{0xF1} (powered, all four channels reported on),
\cs{NR11} to \cs{0xFF} (50\% duty, full length) and \cs{NR51} to
\cs{0xFF} (everything to both sides).  The LFSR starts at \cs{0x7FFF} and
wave RAM at zero.

\subsection{What the APU does not do}
\label{apu-limits}

\begin{itemize}
\item Envelope parameters are not latched from \cs{NRx2} on a plain write
      (see above).
\item Wave RAM is not blocked while channel~3 is enabled, so a test that
      checks the read-back lock will fail.
\item The output level is about half of what a real DMG produces, and the
      \cs{NR50} divide-by-zero case --- both volumes zero --- produces
      silence rather than the 50\%-mute behaviour of real hardware.
\item There is no high-pass filter, so the DC offset that gives the DMG
      its characteristic thin sound is absent.
\item Length-counter reload and trigger timing are modelled in units of
      the frame sequencer rather than at cycle granularity, so the
      cycle-exact \fname{blargg} \fname{dmg_sound} tests still fail even
      where the audible result is right.
\end{itemize}

\attention
The \fname{TODO} file claims that \fname{src/apu.c} is ``minimal (length
counters + NR52 on/off only)''.  That is stale --- the file now contains
the full frame sequencer, sweep, envelope, wave and LFSR synthesis.  Treat
the \fname{TODO} entries about the APU as a record of an older state of
the code, and re-run the suites before acting on them.

\clearpage

% ===========================================================================
\section{Persistence}
\label{persistence}

Two kinds of file are written next to the ROM, both named by appending an
extension to the ROM path.  Neither uses a directory: \fname{gbemu.save}
and \fname{gbemu.state} appear next to the ROM, and if the ROM directory
is not writable both operations fail with a message rather than falling
back to a temporary location.

\subsection{Battery saves}
\label{battery-format}

\fname{gb_battery_save} writes \cs{<rom path>.sav}, but only for
cartridge types that declare a battery-backed RAM:

\begin{longtable}{@{}C{0.16\textwidth}L{0.52\textwidth}C{0.20\textwidth}@{}}
\toprule
\textbf{Type} & \textbf{Cartridge} & \textbf{Payload}\\
\midrule
\endhead
\cs{0x03} & MBC1 + RAM + battery & SRAM\\
\cs{0x06} & MBC2 & 512 bytes of 4-bit RAM\\
\cs{0x09} & ROM + RAM + battery & SRAM\\
\cs{0x0F} & MBC3 + timer + battery & SRAM + 5 RTC bytes\\
\cs{0x10} & MBC3 + timer + RAM + battery & SRAM + 5 RTC bytes\\
\cs{0x13} & MBC3 + RAM + battery & SRAM\\
\cs{0x1B} & MBC5 + RAM + battery & SRAM\\
\cs{0x1E} & MBC5 + rumble + RAM + battery & SRAM\\
other & none & nothing written\\
\bottomrule
\end{longtable}

The SRAM payload is always the full \cs{sram\_size} the loader allocated,
not the size the header asked for --- so the file is the 8\,KiB minimum
for small carts.  For MBC2 it is exactly 512 bytes, the size of
\cs{mbc2\_ram}.

MBC3 timer carts append five bytes in the layout every other emulator
uses: seconds, minutes, hours, the low byte of the day counter, and a
high byte whose bit~0 is the ninth day bit, bit~6 the halt flag and bit~7
the carry flag.  \cs{rtc\_pack} and \cs{rtc\_unpack} are the only places
that conversion exists.

\attention
Because the save covers the whole allocation, a \cs{.sav} written by one
build is only readable by another build of the same emulator, and only for
the same cartridge --- the file has no header.  It is, however, the layout
other emulators use, so it is portable between emulators.

\seealso
\fname{gb_battery_load} reads the same file at start-up, prints
\cs{[INFO] Battery save loaded: <path>} on success, and reports
\cs{[ERR] Battery save is truncated: <path>} if the payload is short.
A missing file is not an error: it simply returns \cs{false} and the game
starts with a fresh save.

\subsection{Save states}

\fname{gb_save_state} writes \cs{<rom path>.state} as a flat
little-endian dump preceded by a 16-byte header:

\begin{Code}
typedef struct {
    char     magic[4];      // "EMU1"
    uint32_t emu_version;   // GB_VERSION_HEX
    uint32_t rom_hash;      // FNV-1a over the padded ROM image
    uint32_t sram_size;     // SRAM allocation at capture time
} StateHeader;
\end{Code}

The payload follows in a fixed order, and the order is part of the file
format:

\begin{enumerate}
\item the whole \cs{gb\_cpu} struct
\item the whole \cs{gb\_timer} struct
\item \cs{gb\_ppu} up to \cs{offsetof(gb\_ppu, frame\_buffer)} --- the
      registers and internal state, but \emph{not} the 92\,KiB of pixels
\item \cs{gb\_apu} up to \cs{offsetof(gb\_apu, audio\_buf)} --- the
      synthesis state, but not the ring buffer or its indices
\item the bus, field by field: the four memory arrays, \cs{io}, \cs{ie},
      the whole DMA state, \cs{double\_speed}, both joypad bytes, the
      mapper type and every mapper's registers (including the full
      \cs{mbc2\_ram}), and finally the SRAM contents
\end{enumerate}

\marginlabel{Field-wise on purpose:}
The bus is written field by field rather than with one
\cs{fwrite(bus, sizeof *bus)} so that the pointers \cs{rom},
\cs{sram}, \cs{timer}, \cs{ppu} and \cs{apu} --- and the
\cs{sram\_dirty} hint --- are automatically excluded.  Two macros,
\cs{W(field)} and \cs{R(field)}, do the work and are \cs{\#undef}'d at the
end of the file.

Two fields are deliberately not saved, and both are recreated or
irrelevant afterwards:

\begin{itemize}
\item The \cs{apu} ring buffer indices are reset to zero after a load, so
      the audio callback cannot consume samples that belonged to the
      moment of capture.
\item \cs{bus->joypad\_interrupt} is not saved.  It is a one-instruction
      edge, and the game will re-raise it if a button is still held.
\end{itemize}

\subsection{Validation}
\label{state-validation}

\fname{gb_load_state} refuses a state that does not match the running
machine, and says why:

\begin{longtable}{@{}L{0.34\textwidth}L{0.58\textwidth}@{}}
\toprule
\textbf{Check} & \textbf{Message on failure}\\
\midrule
\endhead
File exists & \cs{[ERR] No save state found: <path>}\\
Magic is \cs{EMU1} & \cs{[ERR] Not a save state: <path>}\\
\cs{emu\_version == GB\_VERSION\_HEX} &
  \cs{[ERR] Save state gb\_version mismatch: file <n>, emulator <m>}\\
\cs{rom\_hash == fnv1a(rom, rom\_size)} &
  \cs{[ERR] Save state was captured from a different ROM}\\
\cs{sram\_size} matches & \cs{[ERR] Save state SRAM size mismatch: ...}\\
Payload fully readable &
  \cs{[ERR] Save state is truncated: <path>}\\
\bottomrule
\end{longtable}

\marginlabel{Hash choice:}
The ROM fingerprint is FNV-1a, a 32-bit multiply-and-xor with no
multiplication table and no dependencies, computed over the whole
\emph{padded} buffer --- including the zero fill above the file --- and
therefore after the interrupt-vector rewrite of
section~\ref{vector-sanitisation}.  Two builds of the emulator agree only
when \cs{GB\_VERSION\_HEX} agrees, so a state written by a debug build
cannot be loaded by a release build of the same source.

\attention
Save states are raw struct dumps, so they are not portable across
architectures, compilers or struct layouts.  They are a within-session
convenience --- suspend and resume --- not an interchange format.

\subsection{When things are written}

\begin{longtable}{@{}L{0.30\textwidth}L{0.62\textwidth}@{}}
\toprule
\textbf{Moment} & \textbf{Action}\\
\midrule
\endhead
Start-up, before the front end is created &
  \fname{gb_battery_load} for the ROM path\\
\cs{F5} or \cs{START+SELECT} & \fname{gb_save_state}\\
\cs{F9} or \cs{START+SELECT} & \fname{gb_load_state}\\
\fname{gb_loop} returns & \fname{gb_battery_save}, then
  \cs{fe->destroy}\\
\bottomrule
\end{longtable}

A session that ends by killing the process rather than by closing the
window loses both the battery save and any save state made since the last
clean exit.

\clearpage

\clearpage

% ===========================================================================
\section{Front ends}
\label{frontends}

Every front end implements the same four-function vtable and returns a
\cs{gb\_frontend*} from a \cs{frontend\_*\_create} constructor.  Five of
them exist; which one is compiled is decided at build time
(section~\ref{building}) or by the host application.

\begin{longtable}{@{}L{0.26\textwidth}L{0.20\textwidth}L{0.46\textwidth}@{}}
\toprule
\textbf{File} & \textbf{Constructor} & \textbf{Purpose}\\
\midrule
\endhead
\fname{frontend/sdl.c} & \cs{frontend\_sdl\_create(gb\_apu*)} &
  Window, texture, audio device, keyboard.  The default build.\\
\fname{frontend/terminal.c} & \cs{frontend\_terminal\_create()} &
  ANSI true-colour output with half-block scaling, raw-mode
  keyboard.\\
\fname{frontend/headless.c} & \cs{frontend\_headless\_create()} &
  No output, no input, never terminates.  The conformance harness.\\
\fname{frontend/wasm.c} & \cs{frontend\_wasm\_create()} &
  Blits the frame buffer into a canvas through an \cs{EM\_JS} function.\\
\fname{frontend/esp32.c} & \cs{frontend\_esp32\_create(gb\_apu*)} &
  SPI ST7789/ST7735 panel, GPIO buttons, optional I2S DAC, on
  ESP-IDF 5.\\
\bottomrule
\end{longtable}

\subsection{The SDL front end}

\fname{sdl_init} initialises SDL's video and audio subsystems, then
creates a window of \cs{160*scale} by \cs{144*scale} pixels (scale from
\fname{EMU_SCALE}, default 4, clamped to at least 1), a streaming
renderer, and an \cs{SDL\_PIXELFORMAT\_ARGB8888} texture --- which is why
the frame buffer is ARGB and not RGBA.

Audio uses \cs{SDL\_OpenAudioDeviceStream} in
\cs{SDL\_AUDIO\_F32} mono at \cs{GB\_APU\_SAMPLE\_RATE}.  The callback
asks for \cs{additional\_amount} bytes, fills them from
\fname{gb_apu_buf_pop} in chunks of at most 1024 samples, and pushes
them into the stream.  If the emulator has not produced a sample the pop
returns \cs{0.0f}, so a stalled emulator produces silence rather than
stutter.

\fname{sdl_poll_events} drains the queue and maintains the joypad bytes
with read-modify-write operations, so two keys mapped to different bits
never clobber each other.  Both \cs{SDL\_EVENT\_KEY\_DOWN} and
\cs{SDL\_EVENT\_KEY\_UP} are handled, which is why key repeat cannot
accumulate.  \cs{F5} and \cs{F9} fire only on the down transition.

\marginlabel{One quirk:}
\fname{joypad_interrupt} is raised for \emph{any} key press, including
keys that are not mapped to a button.  A game that enables the joypad
interrupt and uses interrupts as a timing source will see extra requests.

\subsection{The terminal front end}
\label{terminal-frontend}

\fname{frontend/terminal.c} needs no graphics library at all.  It renders
by printing one line of terminal rows per \emph{two} Game Boy scanlines,
using the half-block character \cs{U+2580} with a true-colour foreground
for the upper pixel and background for the lower one:

\begin{Code}
for (y = 0; y < height - 1; y += 2)
  for (x = 0; x < width; x++)
    printf("\033[38;2;%d;%d;%dm\033[48;2;%d;%d;%dm\u2580", ...top, ...bottom);
\end{Code}

The row buffer is 8192 bytes, sized for 160 pixels of escape sequences
plus a reset and a newline.  143 rows are emitted, since the last
scanline has no partner.

\fname{term_init} puts the terminal in raw mode
(\cs{ECHO}, \cs{ICANON} and \cs{IEXTEN} cleared, \cs{VMIN} and
\cs{VTIME} zero) so that key presses arrive without a newline, hides the
cursor, clears the screen and disables line wrapping.  It also installs a
\cs{SIGINT} handler that only sets a flag, because the loop must leave
\fname{gb_loop} before the terminal can be restored.
\fname{term_destroy} re-enables the cursor and wrapping and puts the
original \cs{termios} back.

Keys are read with a single non-blocking \cs{read} of up to 64 bytes and
decoded inline:

\begin{longtable}{@{}L{0.30\textwidth}L{0.28\textwidth}L{0.34\textwidth}@{}}
\toprule
\textbf{Input} & \textbf{Button} & \textbf{Note}\\
\midrule
\endhead
\cs{ESC [ A/B/C/D} & Up / Down / Right / Left & Arrow escape
  sequences; the sequence is consumed from the same buffer\\
\cs{z}, \cs{x} & A, B &\\
\cs{0x7F} & Select & The backspace byte\\
\cs{LF} or \cs{CR} & Start &\\
\cs{q}, \cs{Q} & --- & Quit\\
bare \cs{ESC} & --- & Quit\\
\cs{Ctrl-C} & --- & Quit, via the signal flag\\
\bottomrule
\end{longtable}

\attention
Because there is no key-repeat protocol in a raw terminal, each press is
treated as five frames long (\cs{KEY\_HOLD\_FRAMES}), which is about
83\,ms, and released afterwards.  A held arrow key therefore has to be
pressed repeatedly.  For the same reason \cs{term\_poll\_events} raises
\cs{joypad\_interrupt} unconditionally on every poll, so the joypad
interrupt fires sixty times a second whether or not anything was pressed.

\subsection{The headless front end}

Three functions that do nothing and one that frees a \fname{NULL}.  It
never clears the \cs{running} flag, so a headless process runs until it is
killed or signalled.  That is exactly what \fname{scripts/test} depends
on: each ROM is run under \cs{gtimeout} and the harness reads its standard
output.  Any front end that returned \cs{false} from \cs{poll\_events}
would end the run immediately and break the harness.

\subsection{The WebAssembly front end}

\fname{frontend/wasm.c} contains one \cs{EM\_JS} function that copies the
frame buffer into a \cs{ImageData} and paints it:

\begin{Code}
EM_JS(void, js_render_to_canvas, (const void* buffer, int w, int h), {
    var buf32 = new Uint32Array(Module.HEAPU8.buffer, buffer, w * h);
    for (var i = 0; i < w * h; i++) {
        data[i*4]     = (buf32[i] >> 16) & 0xFF;   // ARGB -> RGBA
        data[i*4 + 3] = 0xFF;
    }
    ctx.putImageData(imageData, 0, 0);
});
\end{Code}

The red and green channels are assigned inside the loop but are always
\cs{0} for the current palette, since the emulator writes
\cs{0xFF000000} black and \cs{0xFFFFFFFF} white.  \cs{init} only sets the
canvas element's \cs{width} and \cs{height} attributes; \cs{poll\_events}
is empty, because input arrives from JavaScript (see
section~\ref{wasm-host}).

\subsection{The ESP32 front end}
\label{esp32-frontend}

\fname{frontend/esp32.c} is not built by \cs{make}; it is meant to be
compiled inside an external ESP-IDF (5.0 or newer) project together with
the core.  Every piece of hardware configuration is a macro with a
default, so the file can be dropped into a project for a different board
without edits:

\begin{longtable}{@{}L{0.30\textwidth}C{0.62\textwidth}@{}}
\toprule
\textbf{Group} & \textbf{Macros (defaults)}\\
\midrule
\endhead
SPI & \cs{EMU\_LCD\_SPI\_HOST} (\cs{SPI2\_HOST}), \cs{EMU\_LCD\_PIN\_SCLK}
  (18), \cs{MOSI} (23), \cs{CS} (15), \cs{DC} (2), \cs{RST} (4),
  \cs{BL} (21), \cs{EMU\_LCD\_BL\_ACTIVE\_LOW} (0)\\
Panel & \cs{EMU\_LCD\_PANEL\_ST7789} (1), \cs{EMU\_LCD\_WIDTH} (240),
  \cs{HEIGHT} (240), \cs{EMU\_LCD\_OFF\_X}/\cs{Y} (centred),
  \cs{EMU\_LCD\_SPI\_CLOCK\_HZ} (80\,MHz), \cs{MIRROR\_X}/\cs{Y}
  (\cs{false}), \cs{COLOR\_INVERT} (0)\\
Buttons & \cs{EMU\_BTN\_UP} (32), \cs{DOWN} (33), \cs{LEFT} (25),
  \cs{RIGHT} (26), \cs{A} (27), \cs{B} (14), \cs{SELECT} (13),
  \cs{START} (19)\\
Hotkeys & \cs{EMU\_HOTKEY\_SAVE\_FRAMES} (30), \cs{LOAD\_FRAMES} (150)\\
Audio & \cs{EMU\_ENABLE\_AUDIO} (1), \cs{EMU\_I2S\_PIN\_BCLK} (27),
  \cs{LRCK} (26), \cs{DOUT} (25), \cs{EMU\_AUDIO\_TASK\_CORE} (1)\\
\bottomrule
\end{longtable}

Rendering converts ARGB8888 to RGB565 with the bytes swapped, because SPI
panels shift out the most significant byte first, and blits with
\cs{esp\_lcd\_panel\_draw\_bitmap} into the DMA-capable buffer.  The frame
is centred in the panel using the offset macros.

Audio runs in its own pinned FreeRTOS task which pops
\fname{gb_apu_buf_pop} in blocks of 256 samples, converts to 16-bit
integers and writes them to an I2S standard-mode channel; the blocking
write is what paces the task.  If either the channel allocation or the
initialisation fails, audio is disabled and the emulator still runs.

\attention
The default GPIO assignments for the buttons and the I2S pins overlap ---
\cs{EMU\_BTN\_B} is 14 while \cs{EMU\_I2S\_PIN\_DOUT} is 25 and
\cs{EMU\_BTN\_A} is 27 while \cs{EMU\_I2S\_PIN\_BCLK} is 27 --- so a board
that enables audio \emph{and} the A/B buttons on the default pins will
not work.  Override at least the I2S pins.  The \cs{START} default of 19 is
deliberately not 16 or 17, which are taken by PSRAM on WROVER boards.

\seealso
The \cs{START + SELECT} combination is the only input on a device with no
keyboard: 30 consecutive frames (about half a second) saves, 150 frames
(about two and a half seconds) loads.  Releasing either button resets the
counter.

\clearpage

% ===========================================================================
\section{Embedding the core}
\label{embedding}

The library exposes everything an application needs and nothing else; the
question is only which parts of the start-up sequence an embedding
application is responsible for.

\subsection{The minimal session}

This is the whole of a working program apart from the front end:

\begin{Code}
#include "gbemu.h"

gb_cpu   cpu;
gb_bus   bus;
gb_timer timer;
gb_ppu   ppu;
gb_apu   apu;

bus.timer = &timer;   // wire the bus to its components FIRST
bus.ppu   = &ppu;
bus.apu   = &apu;

gb_cpu_init(&cpu);
gb_timer_init(&timer);
gb_ppu_init(&ppu);
gb_apu_init(&apu);

if (!gb_bus_load_rom(&bus, "game.gb")) return 1;
gb_battery_load(&bus, "game.gb");

bus.io[0x0F] = 0x01;     // IF: V-blank request left by the boot ROM
bus.ie       = 0x00;
bus.joypad_dpad    = 0x0F;   // all released
bus.joypad_buttons = 0x0F;

gb_frontend* fe = my_frontend_create();
fe->init(fe, GB_SCREEN_WIDTH, GB_SCREEN_HEIGHT);
gb_loop(&cpu, &bus, &timer, &ppu, &apu, fe);
gb_battery_save(&bus, "game.gb");
fe->destroy(fe);
\end{Code}

\attention
Four details are not optional:

\begin{enumerate}
\item \textbf{Wire the pointers first.}  \cs{bus->ppu} and friends are
      tested for \fname{NULL} inside the bus, so a missing one is a silent
      no-op rather than a crash.
\item \textbf{Call the four \cs{init} functions.}  \cs{gb\_cpu\_init} sets
      \cs{PC} and the registers, \cs{gb\_ppu\_init} the video registers
      and clears the frame buffer, \cs{gb\_apu\_init} the sound defaults,
      and \cs{gb\_timer\_init} seeds the divider with \cs{0xABCC}.
\item \textbf{Set the post-boot interrupt state.}  Neither
      \cs{gb\_cpu\_init} nor \cs{gb\_timer\_init} touches \cs{IF},
      \cs{IE} or the joypad bytes; the host has to.
\item \textbf{Decide \cs{A} from the cartridge.}  Both
      \fname{apps/cli.c} and \fname{apps/wasm.c} force
      \cs{cpu.a = 0x01} when header byte \cs{0x0143} is neither \cs{0x80}
      nor \cs{0xC0}, because \cs{gb\_cpu\_init} leaves the CGB boot value
      \cs{0x11} in place.  Skipping this leaves DMG games in CGB mode.
\end{enumerate}

\subsection{Running without a loop}

\fname{gb_loop} is a convenience, not a requirement.  It is a plain
frame-budgeted loop and can be replaced entirely; the WASM host already
does exactly that (section~\ref{wasm-host}).  The minimum sequence for
one frame is:

\begin{Code}
int dots = 0;
while (dots < GB_CYCLES_PER_FRAME) {
    int t = gb_cpu_step(&cpu, &bus);
    int scale = bus.double_speed ? 1 : 2;
    dots += t * scale;
    int it = gb_handle_interrupts(&cpu, &bus, &ppu, &timer);
    dots += it * scale;
}
if (ppu.frame_ready) { /* draw */ ppu.frame_ready = false; }
\end{Code}

Any front-end-like policy --- rendering, input, sleeping, audio
draining --- can then be woven in wherever it suits the host.

\subsection{Hotkeys and save states}
\label{embedding-hotkeys}

The application, not the front end, decides what a hotkey means.  Both
\fname{apps/cli.c} and an embedding application pass a context pointer
and a callback:

\begin{Code}
typedef struct {
    gb_cpu* cpu; gb_bus* bus; gb_timer* timer;
    gb_ppu* ppu; gb_apu* apu; const char* rom_path;
} EmuContext;

static void on_hotkey(void* userdata, gb_hotkey key) {
    EmuContext* ctx = userdata;
    switch (key) {
    case GB_HOTKEY_SAVE_STATE:
        gb_save_state(ctx->cpu, ctx->bus, ctx->timer, ctx->ppu, ctx->apu,
                      ctx->rom_path);
        break;
    case GB_HOTKEY_LOAD_STATE:
        gb_load_state(ctx->cpu, ctx->bus, ctx->timer, ctx->ppu, ctx->apu,
                      ctx->rom_path);
        break;
    }
}

fe->hotkey_ctx = &ctx;
fe->on_hotkey  = on_hotkey;
\end{Code}

A front end that has no use for them leaves \cs{on\_hotkey} \fname{NULL},
which is what the terminal, headless and WASM front ends do.

\subsection{Audio consumer}

The APU never blocks and never calls into the host.  It produces samples
at \cs{GB\_APU\_SAMPLE\_RATE} and drops them if nobody is reading, so the
only requirement on a host is to call \fname{gb_apu_buf_pop}
periodically:

\begin{Code}
float s = gb_apu_buf_pop(&apu);   // 0.0f when the buffer is empty
\end{Code}

It is safe to do this from another thread --- the ring buffer is
single-producer single-consumer --- and both real front ends do.  Nothing
else in the core is thread-safe: \cs{gb\_cpu\_step},
\cs{gb\_bus\_read}, \fname{gb_bus_write} and \cs{gb\_save\_state} must all
be called from the emulation thread.

\subsection{Linking}

\begin{Code}
cc -Isrc my_game.c -L. -lgbemu $(pkg-config --cflags --libs sdl3) -o my_game
cc -Isrc -DGBEMU_STRIP_PREFIX my_game.c libgbemu.a -o my_game
\end{Code}

The static archive has the front ends \emph{out} of it --- only
\fname{src/*.c} is archived --- so an embedding application links its own
front end, or none at all.  That is why the three constructor prototypes
in \fname{apps/cli.c} are declared locally rather than included from a
header: exactly one of them is defined in the link.

\attention
There is no \cs{main} in the archive.  \fname{apps/cli.c} and
\fname{apps/wasm.c} are the only programs, and both are compiled directly
rather than archived.

\subsection{The WebAssembly host}
\label{wasm-host}

The browser build inverts the usual structure: the C side is driven from
JavaScript and has no input loop of its own.

\begin{longtable}{@{}L{0.28\textwidth}L{0.62\textwidth}@{}}
\toprule
\textbf{Exported symbol} & \textbf{Meaning}\\
\midrule
\endhead
\cs{wasm\_load\_rom} &
  Zeroes all five state objects, re-runs the \cs{init} functions and
  loads \cs{/rom.gb} from the virtual file system.  Safe to call again
  to reset the session.\\
\cs{wasm\_set\_key(key, pressed)} &
  Maps a browser \cs{keyCode} to a joypad bit: 37--40 for the D-pad,
  90 (\cs{Z}) for A, 88 (\cs{X}) for B, 8 for Select and 13 for
  \cs{Enter}.  Raises the joypad interrupt on press.\\
\cs{wasm\_is\_running} &
  Returns 1 while a ROM is loaded and 0 otherwise, so the shell can
  decide whether to show the picker.\\
\cs{main} &
  Creates the front end, loads \cs{/rom.gb} if it exists, and installs
  the main loop with \cs{emscripten\_set\_main\_loop} in
  unlimited-render mode.\\
\bottomrule
\end{longtable}

\fname{web/shell.html} is the shell: a canvas with
\cs{id="screen"}, a ROM picker, and the keyboard handling that calls
\cs{wasm\_set\_key}.  Because the loop runs in the browser's
\cs{requestAnimationFrame}, frame pacing is the browser's job and the
\cs{EMU\_NOSLEEP} logic is not compiled in.

\clearpage

\clearpage

% ===========================================================================
\section{Testing}
\label{testing}

Correctness is measured against three external ROM suites that are run
against a headless build of the emulator.  No unit tests exist: the
project has one test entry point, \fname{scripts/test}, and its
assertions live inside the ROMs.

\subsection{The suites}

\begin{longtable}{@{}L{0.22\textwidth}L{0.30\textwidth}L{0.40\textwidth}@{}}
\toprule
\textbf{Suite} & \textbf{Source} & \textbf{What it checks}\\
\midrule
\endhead
\fname{gb} & \cs{tests/gb-test-roms} &
  Blargg's classic test ROMs, both as composites and as individual
  cases: \cs{instr\_timing}, \cs{interrupt\_time}, \cs{halt\_bug},
  \cs{cpu\_instrs}, \cs{mem\_timing}, \cs{mem\_timing-2},
  \cs{oam\_bug}, \cs{dmg\_sound}, \cs{cgb\_sound}\\
\fname{samesuite} & \cs{tests/SameSuite} &
  SameBoy's hardware-behaviour ROMs, one per behaviour: \cs{cpu},
  \cs{dma}, \cs{ppu}, \cs{apu}, \cs{timer}, \cs{interrupt}, \cs{sgb} and
  so on\\
\fname{mooneye} & \cs{tests/mooneye\-test\-suite} &
  The acceptance and emulator-only ROMs whose hardware suffix says they
  are expected to pass on a DMG\\
\bottomrule
\end{longtable}

SameSuite ROMs are assembled with \textbf{rgbds} and mooneye ROMs with
\cs{wla\-dx}; the script builds them if the tools are present, and skips
the build if prebuilt ROMs already exist.

\marginlabel{mooneye filtering:}
mooneye names carry a hardware suffix after the last hyphen --- an empty
suffix means ``all models'', \cs{-GS} means Game Boy and Super Game Boy,
\cs{-dmgABCmgb} names specific models, \cs{-sgb} is Super Game Boy only.
The harness keeps a test only if there is no suffix, or the suffix
contains \cs{dmg} or a capital \cs{G}.  This is why a CGB-only test does
not appear as a failure: it is never run.

\subsection{Verdicts}

Verdicts arrive over the emulated serial port, using the two interception
paths of section~\ref{serial}.

\begin{description}

\item[SameSuite and mooneye] The ROM writes a result byte per iteration
      through \cs{SB}/\cs{SC}.  Success is the Fibonacci register
      \cs{03 05 08 0D 15 22}; failure is six \cs{0x42} bytes.  The script
      pipes the emulator's standard output through \cs{xxd -p} and
      compares the hex string.
\item[Blargg] The ROM prints human-readable lines through cartridge RAM
      at \cs{0xA004} and finishes with \cs{Passed} or \cs{Failed}.  The
      script greps for \cs{passed} while requiring that \cs{failed} does
      \emph{not} also appear, which distinguishes a partial composite
      report from a clean pass.
\end{description}

\attention
The \cs{Passed} test runs before the \cs{Failed} test on purpose: the
bl'rg composites print every individual result and a final verdict, and a
ROM that prints both \cs{01:ok 02:ok ... Failed} must be scored as a
failure.

\subsection{Running the harness}

\begin{Code}
./scripts/test                # all three suites
./scripts/test gb             # only the bl'rg ROMs
./scripts/test samesuite      # only SameSuite
./scripts/test mooneye        # only mooneye
./scripts/test --help
\end{Code}

The script sets two environment variables that matter:
\cs{EMU\_NOSLEEP=1} so that each ROM runs as fast as the CPU allows, and
\cs{EMU\_HEADLESS=1} so that no window is created.  It then rebuilds the
emulator:

\begin{Code}
make clean
make all type=RELEASE HEADLESS=1
\end{Code}

\attention
The rebuild is unconditional and it \emph{overwrites} a debug build.  If
you are running the harness under a debugger, or have
\fname{libgbemu.a} linked into something else, build the emulator again
afterwards --- and remember that the release build has no sanitizers, so
memory errors that the debug build would report here simply corrupt the
picture.

Each ROM runs under \cs{gtimeout} with a suite-specific limit: 60 seconds
for the bl'rg composites, 30 seconds for individual bl'rg ROMs, and 15
seconds for SameSuite and mooneye, which print their result within a few
seconds and then halt forever.

\subsection{Report files}

Three reports are written, one per suite:
\fname{test-report-gb.txt}, \fname{test-report-samesuite.txt} and
\fname{test-report-mooneye.txt}.  Each test contributes one block:

\begin{Code}
TEST: [composite] instr_timing
  File:  .../tests/gb-test-roms/instr_timing/instr_timing.gb
  Lines: 4
  Status: PASS
  Detail: instr_timing; ; ; Passed;
---
\end{Code}

\cs{Status} is one of six values:

\begin{longtable}{@{}L{0.18\textwidth}L{0.74\textwidth}@{}}
\toprule
\textbf{Status} & \textbf{Meaning}\\
\midrule
\endhead
\cs{PASS} & The expected signature was seen\\
\cs{FAIL} & The failure signature was seen, or the text verdict said
  \cs{Failed}, or nothing was printed at all (\cs{NOOUT})\\
\cs{NOOUT} & The emulator exited without producing any output\\
\cs{TIMEOUT} & The emulator was killed having produced nothing\\
\cs{PARTIAL} & The emulator was killed having produced a prefix of the
  signature --- the most useful status, since it says how far it got\\
\cs{UNKNOWN} & Some other output\\
\bottomrule
\end{longtable}

A per-suite summary with the counts of each status is appended to each
report.

\subsection{Reading a failure}

\begin{enumerate}
\item Look at \cs{Detail}.  For a SameSuite or mooneye ROM it is the raw
      serial bytes; \cs{4242} means the ROM decided it had failed.
\item A \cs{PARTIAL} serial result is not a failure of the emulator --- it
      means the ROM printed its result and then kept running, so it
      usually indicates that the test itself passed and the timeout was
      simply too short, or that the ROM halted before finishing its
      remaining iterations.
\item Check \cs{Lines} for the bl'rg composites: a composite that dies
      early prints far fewer lines than the last one in a passing run.
\item Cross-check the group name against
      section~\ref{known-gaps}.  Most failures have a documented cause,
      and re-running the suite after fixing one of them usually moves the
      failure elsewhere rather than removing it.
\end{enumerate}

\attention
\cs{/tmp/blargg\_raw.bin} accumulates every byte written to cartridge RAM
at \cs{0xA004} and above, including non-printable ones, for the whole
process.  It is truncated when the first write happens, so run one ROM at a
time if you want to inspect it.  The file is not created by the ESP32
build.

\subsection{What the suites do not cover}

\begin{itemize}
\item No multi-frame determinism test: there is no test that runs the same
      ROM twice and compares the frame buffers.
\item No fuzzing of the bus, so an access to an address the mapper model
      does not expect is caught only if a ROM happens to make it.
\item No performance regression tracking.  \cs{EMU\_NOSLEEP=1} makes the
      harness run as fast as possible but nothing measures the rate.
\item The bl'rg and mooneye suites are third-party; \cs{make all} does not
      fetch them, and the script fails with an explicit message if a
      suite's ROMs are neither present nor buildable.
\end{itemize}

\clearpage

% ===========================================================================
\section{Extending the emulator}
\label{extending}

This chapter is about changing the code without breaking its two
contracts: the cycle count an instruction reports, and the size and
layout of the structs a save state dumps.

\subsection{Adding an instruction}

The generated table has to know about it first.  Add the entry to
\fname{docs/dev/Opcodes.json} with its mnemonic, operands, size and cycle
counts, then regenerate:

\begin{Code}
python3 scripts/generate_opcodes.py
\end{Code}

The script reads the JSON, derives the \cs{OP\_*} enum name from the
mnemonic and operands using the rules in
section~\ref{opcode-tables}, and rewrites \fname{src/opcodes.h} from
scratch.  Editing \fname{opcodes.h} by hand is pointless: the next
regeneration discards it.

Then implement the body in the handler that already owns that category,
and add the \cs{case} label:

\begin{Code}
case OP_SLA_A: {
    uint8_t carry = (cpu->a >> 7) & 1;
    cpu->a = cpu->a << 1;
    gb_flag_set(cpu, GB_FLAG_Z, cpu->a == 0);
    gb_flag_set(cpu, GB_FLAG_N, false);
    gb_flag_set(cpu, GB_FLAG_H, false);
    gb_flag_set(cpu, GB_FLAG_C, carry);
    break;
}
\end{Code}

\attention
Three things must all stay true, and nothing checks any of them:

\begin{enumerate}
\item \textbf{The number of machine cycles matches the table.}  Count the
      \cs{gb\_machine\_tick(bus, 4)} calls made on the execution path,
      including the ones inside \fname{gb_fetch8}, \fname{gb_fetch16} and
      the register-index helpers, and add the internal cycles.  The total
      must equal the entry's cycle count divided by four.
\item \textbf{Flags are set explicitly.}  Every arithmetic or logical
      instruction must write \cs{Z}, \cs{N} and \cs{H} even when the answer
      is ``unchanged'': on hardware those flags are not preserved.
\item \textbf{\cs{POP AF} masks the low nibble.}  If your instruction can
      affect \fname{F} indirectly, check that path.
\end{enumerate}

The bl'rg \cs{cpu\_instrs} and SameSuite \cs{cpu} suites are the tests that
will find a mistake here.

\subsection{Adding a component}

A new emulated device is a struct plus a \cs{step} function plus a call
site in \fname{gb_machine_tick}:

\begin{Code}
typedef struct gb_something { /* ... */ } gb_something;

void gb_something_init(gb_something* d);
void gb_something_step(gb_something* d, int dots);
uint8_t gb_something_read(const gb_something* d, uint16_t addr);
void    gb_something_write(gb_something* d, uint16_t addr, uint8_t v);

// in gb_machine_tick
if (bus->something) gb_something_step(bus->something, t_cycles * scale);
\end{Code}

\attention
If you also want it in save states, add it to the \cs{W} and \cs{R} chains
in \fname{save.c} --- in the same position on both sides, since the format
is positional.  And keep the large buffers at the \emph{end} of the
struct, because \fname{gb_save_state} serialises a component with
\fname{fwrite(ppu, offsetof(gb_ppu, frame_buffer), 1, f)} and a new member
inserted before the buffer would silently be included.

Bus routing follows the existing pattern: test
\cs{bus->something != NULL} before delegating, and fall through to the
\cs{io} shadow array otherwise, so a bus without the component still
answers reads.

\subsection{Adding a front end}

Write four functions and one constructor, following
\fname{frontend/headless.c}, which is the whole of a working (if useless)
front end:

\begin{Code}
static bool  my_init  (gb_frontend* fe, int w, int h);
static void  my_render(gb_frontend* fe, const uint32_t* buf, int w, int h);
static void  my_poll  (gb_frontend* fe, gb_bus* bus, bool* running);
static void  my_destroy(gb_frontend* fe);

gb_frontend* my_frontend_create(void) {
    gb_frontend* fe = calloc(1, sizeof(gb_frontend));
    fe->priv = my_private_state();
    fe->init = my_init;
    fe->render = my_render;
    fe->poll_events = my_poll;
    fe->destroy = my_destroy;
    return fe;
}
\end{Code}

Rules that the loop depends on:

\begin{itemize}
\item \cs{init} must return \cs{false} rather than aborting; \cs{main}
      turns that into exit status~1.
\item \cs{render} must not keep the buffer pointer.  The core clears
      \cs{frame\_ready} immediately after the call and the pixels are
      overwritten on the next scanline.
\item \cs{poll\_events} is the only way to stop the machine, through the
      \cs{running} out-parameter.  A front end that never clears it can
      only be terminated from outside --- which is the headless design.
\item \cs{destroy} must release what \cs{init} allocated.  The core does
      not touch the vtable again.
\end{itemize}

Then add it to \fname{make/common.mk} as another
\cs{FRONTEND\_SRC} case, and add the corresponding \cs{-D} test to
\fname{apps/cli.c}.

\subsection{Adding an interrupt}

\begin{enumerate}
\item Reserve the bit in \fname{gb_bus}: a boolean next to the existing
      request flags is enough.
\item Set it where the condition is detected, e.g.\ at the end of a
      scanline in \fname{gb_ppu_step}.
\item Latch it into \cs{IF} in \fname{gb_handle_interrupts}.  That
      function is the \emph{only} place hardware flags become register
      bits, and the pending mask there is
      \cs{if\_reg \& bus->ie \& 0x1F} --- a new bit needs widening, or it
      will be invisible.
\item Add the vector case to the loop in the same function.  The dispatch
      is a linear scan from bit~0, so the loop bound and the vector
      arithmetic (\cs{0x0040 + 8*bit}) both need updating.
\item Extend the save state if the flag is persistent.
\end{enumerate}

\attention
The serial interrupt (\cs{IF} bit~3) is the worked example of how not to
do this: \cs{IF} bit~3 is reachable and dispatchable, but nothing ever
raises it, because serial output is handled by interception instead.  The
vector at \cs{0x0058} exists and costs nothing until it is used.

\subsection{Changing the timing model}

Two constants govern everything, and both are derived from the clock
rather than hard-coded:

\begin{longtable}{@{}L{0.32\textwidth}L{0.58\textwidth}@{}}
\toprule
\textbf{Macro} & \textbf{Definition}\\
\midrule
\endhead
\cs{GB\_CPU\_FREQ} & \cs{4194304} T-cycles per second, the DMG rate\\
\cs{GB\_CYCLES\_PER\_FRAME} & \cs{GB\_CPU\_FREQ / 60}, i.e.\ 69\,905
  T-cycles\\
\cs{GB\_FRAME\_TIME\_NS} & \cs{1000000000L / 60}\\
\bottomrule
\end{longtable}

Note that \cs{GB\_CYCLES\_PER\_FRAME} is expressed in T-cycles but is
compared in \fname{gb_loop} against a running total of \emph{dots}, so
the two agree only because one T-cycle is two dots at normal speed.  That
is not an accident: it is what keeps the frame budget correct in CGB
double-speed mode, where the scale flips to 1.  Any change to the cycle
handling has to preserve that relationship.

\clearpage

% ===========================================================================
\section{Known gaps}
\label{known-gaps}

This chapter collects, in one place, the hardware behaviour the emulator
does not implement.  Each entry says what is missing and which suite
notices.  Most of them were found by running the suites, not by reading
the Game Boy documentation.

\subsection{CPU and timing}

\begin{longtable}{@{}L{0.30\textwidth}L{0.60\textwidth}@{}}
\toprule
\textbf{Gap} & \textbf{Effect}\\
\midrule
\endhead
Atomic instructions &
  An instruction executes as a unit: its total cost is returned from the
  table, and its immediate reads and writes happen in handler order
  rather than at fixed dots.  The bl'rg \cs{mem\_timing} suites, which
  read a register at a precise dot inside \cs{LDH A,(a8)}, fail all
  their cases.  See also
  \cs{ADD SP,e8}/\cs{LD HL,SP+e8}, whose internal cycles are modelled as
  plain increments.\\
\cs{EI} and \cs{HALT} &
  The interaction is not exact: \cs{ime\_scheduled} is promoted at the
  top of \fname{gb_cpu_step}, before the halted check, so a
  \cs{EI;HALT} sequence does not wake the CPU the way hardware does.
  SameSuite \cs{interrupt/ei\_delay\_halt} fails; the bl'rg
  \cs{halt\_bug} composite passes.\\
\cs{STOP} &
  Performs the CGB speed switch and consumes its operand byte, then
  returns.  It does not stop the CPU, halt the divider, mute the APU or
  reset the joypad matrix, all of which real hardware does.\\
\cs{DAA} &
  Implemented, but the C-64 (non-carry) case only ever clears \cs{C} when
  the addition overflowed, rather than preserving it.\\
\bottomrule
\end{longtable}

\subsection{Memory, DMA and the bus}
\label{dma-limits}

\begin{longtable}{@{}L{0.30\textwidth}L{0.60\textwidth}@{}}
\toprule
\textbf{Gap} & \textbf{Effect}\\
\midrule
\endhead
No DMA register validation &
  The OAM DMA source page is not masked and not checked against
  \cs{0xXX00--0xXX9F}.  A page of \cs{0xFE} copies \cs{0xFF} bytes,
  because OAM reads are themselves blocked during the transfer.\\
No DMA restart timing edge &
  A \cs{0xFF46} write during the initial two-cycle delay, or while a
  restart is already pending, restarts immediately instead of after two
  cycles (section~\ref{oam-dma}).\\
No CGB DMA &
  \cs{GDMA} and \cs{HDMA} at \cs{0xFF51}--\cs{0xFF55} are not implemented
  at all, so CGB games that use them hang.  SameSuite \cs{dma/hdma\_lcd\_off}
  and friends fail.\\
Unmapped I/O reads \index{unmapped I/O} &
  Fall through to the \cs{io} shadow array, which reads back whatever was
  last written --- zero initially --- instead of the \cs{0xFF} real
  hardware returns.\\
Echo RAM at \cs{0xFD00}--\cs{0xFDFF} &
  Correctly aliases \cs{0xC000}--\cs{0xCFFF}, so games that probe
  \cs{0xFFFE} in the echo region see HRAM contents rather than a fault.\\
\bottomrule
\end{longtable}

\subsection{Mappers}

\begin{longtable}{@{}L{0.30\textwidth}L{0.60\textwidth}@{}}
\toprule
\textbf{Gap} & \textbf{Effect}\\
\midrule
\endhead
MBC3 RTC is not driven by a clock &
  The latch sequence advances a state machine but never copies the host
  time into the registers, and the counters only advance when a game
  writes them.  A game that reads the clock sees whatever it found at
  load time.  There is a \cs{TODO} marker for this.\\
Serial text output for MBC2 and MBC3 &
  The \cs{0xA004} interception is after the MBC2 and MBC3 write handlers,
  which return early, so text diagnostics from those mappers never reach
  standard output.  Only the \cs{SB}/\cs{SC} path works for them.\\
MBC3 does not distinguish RAM banks 0--3 by size &
  The bank mask is \cs{ram\_bank \& (sram\_banks - 1)}, so a cart
  declaring 8\,KiB of RAM has only one reachable bank and selecting
  bank~3 silently wraps.\\
No MBC6, MBC7, HuC1, camera or rumble &
  Cartridge types outside the four supported families are treated as
  ROM-only.  A rumble cartridge runs, but the motor bit at
  \cs{0x6000} is discarded.\\
\bottomrule
\end{longtable}

\subsection{PPU}

See section~\ref{ppu-limits} for the full list.  In brief: no per-pixel
timing, no VRAM/OAM access blocking, no \cs{STAT} blocking flag, no
\cs{LY}-write side effects, no DMG OAM scan bug, sprite BG priority not
applied, no CGB palette RAM or tile attributes, no VRAM banking.

\subsection{APU}

See section~\ref{apu-limits}.  In brief: envelope parameters not latched
on write, no wave RAM locking, no high-pass filter, output level about
half scale, and cycle-exact length/trigger timing is not modelled.

\subsection{CGB}
\label{cgb-limits}

\begin{longtable}{@{}L{0.30\textwidth}L{0.60\textwidth}@{}}
\toprule
\textbf{Missing} & \textbf{Consequence}\\
\midrule
\endhead
Palette RAM (\cs{BGPI}/\cs{BGPD}) &
  Colour~0 of every tile uses the DMG greyscale table, so CGB games
  display in greyscale.  SameSuite
  \cs{ppu/blocking\_bgpi\_increase} fails.
\end{longtable}

Everything else about the CGB --- VRAM banking, tile attributes, the
separate window line counter, \cs{OPRB}, HDMA, the boot ROM and the
\cs{KEY1} management registers beyond bit~0 --- is absent.
\cs{double\_speed} and \cs{STOP} do work, so a CGB game that only checks
which mode it is in will behave correctly.

\subsection{SGB}

Nothing is implemented: no \cs{MLT\_REQ}, no border palette, no attribute
files, no joypad comms.  Two SameSuite tests,
\cs{sgb/command\_mlt\_req} and \cs{sgb/command\_mlt\_req\_1\_incrementing},
fail; the \fname{TODO} file suggests either implementing the commands or
filtering the tests out.

\subsection{Front ends}

\begin{itemize}
\item The SDL and terminal front ends raise \cs{joypad\_interrupt} more
      eagerly than hardware: SDL does so for any key, the terminal for
      every poll.
\item The terminal front end has no key repeat, so buttons are five
      frames long and must be re-pressed.
\item The headless front end provides no way to end a session, and no
      way to observe anything.
\item The ESP32 front end's default pin assignments conflict between the
      buttons and the I2S pins (section~\ref{esp32-frontend}).
\end{itemize}

\subsection{Stale documentation}

\attention
The \fname{TODO} file at the repository root is out of date.  Among other
things it claims that \fname{src/apu.c} contains only length counters and
that \fname{src/bus.c} has no DMA support --- both of which have since
been implemented.  Read it as a record of what was missing when the entry
was written, and confirm against the source before acting on an item.

\clearpage

% ===========================================================================
\appendix
\label{appendices}

% ===========================================================================
\section{Opcode coverage}
\label{app-coverage}

Of the 256 base opcode values, 244 are dispatched, one is the \cs{CB}
prefix, and eleven are illegal and fatal:

\begin{longtable}{@{}C{0.24\textwidth}L{0.66\textwidth}@{}}
\toprule
\textbf{Opcode} & \textbf{Status}\\
\midrule
\endhead
\cs{0xCB} & The prefix byte; handled by \fname{instr_cb}, never by a
  handler\\
\cs{0xD3}, \cs{0xDB}, \cs{0xDD} & Illegal\\
\cs{0xE3}, \cs{0xE4} & Illegal\\
\cs{0xEB}, \cs{0xEC}, \cs{0xED} & Illegal\\
\cs{0xF4} & Illegal\\
\cs{0xFC}, \cs{0xFD} & Illegal\\
\bottomrule
\end{longtable}

The generated table prices all of them, including the illegal ones, so
\fname{get_opcode_cycles} has no hole in it.  Reaching any of them
prints \cs{Unhandled opcode: 0xNN at PC: 0xNNNN} and calls \cs{exit(1)}.

The eleven handlers own the remaining 244:

\begin{longtable}{@{}L{0.32\textwidth}C{0.16\textwidth}C{0.42\textwidth}@{}}
\toprule
\textbf{Handler} & \textbf{Opcodes} & \textbf{Category}\\
\midrule
\endhead
\fname{instr_8bit_arithm} & 61 & 8-bit arithmetic and logic\\
\fname{instr_load} & 92 & Loads and stack-relative loads\\
\fname{instr_16bit_arithm} & 13 & 16-bit arithmetic\\
\fname{instr_bitwise} & 28 & Bitwise and compare\\
\fname{instr_bit_flag} & 0 & Empty placeholder\\
\fname{instr_bit_shift} & 4 & Accumulator rotates\\
\fname{instr_jumps} & 30 & Jumps, calls, returns\\
\fname{instr_carry} & 2 & Carry flag control\\
\fname{instr_interrupts} & 3 & \cs{DI}, \cs{EI}, \cs{HALT}\\
\fname{instr_stack} & 8 & \cs{PUSH} and \cs{POP}\\
\fname{instr_misc} & 3 & \cs{DAA}, \cs{NOP}, \cs{STOP}\\
\midrule
\textbf{Total} & \textbf{244} &\\
\bottomrule
\end{longtable}

The \cs{CB} space is dispatched by field rather than by table, so all 256
values are covered: 8 rotates/shifts, 8 \cs{BIT}, 8 \cs{RES} and 8
\cs{SET} per register, including \cs{(HL)}.

\clearpage

% ===========================================================================
\section{I/O register map}
\label{app-registers}

The registers the emulator decodes, with their destinations.  Everything
else in \cs{0xFF00--0xFF7F} is a plain byte in \cs{bus->io}.

\begin{longtable}{@{}C{0.16\textwidth}L{0.28\textwidth}L{0.46\textwidth}@{}}
\toprule
\textbf{Address} & \textbf{Register} & \textbf{Handled by}\\
\midrule
\endhead
\cs{0xFF00} & \cs{P1} & \cs{gb\_bus\_read}/\cs{write}, joypad matrix\\
\cs{0xFF01} & \cs{SB} & \cs{io} shadow; source for the serial print\\
\cs{0xFF02} & \cs{SC} & Special: writing \cs{0x81} prints \cs{SB}\\
\cs{0xFF04} & \cs{DIV} & \cs{gb\_timer\_read}/\cs{write}\\
\cs{0xFF05} & \cs{TIMA} & \cs{gb\_timer\_read}/\cs{write}\\
\cs{0xFF06} & \cs{TMA} & \cs{gb\_timer\_read}/\cs{write}\\
\cs{0xFF07} & \cs{TAC} & \cs{gb\_timer\_read}/\cs{write}\\
\cs{0xFF0F} & \cs{IF} & Special: reads OR \cs{0xE0}, writes mask
  \cs{0x1F}\\
\cs{0xFF10--0xFF2F} & \cs{NR10}--\cs{NR51} & \cs{gb\_apu\_read}/\cs{write}\\
\cs{0xFF30--0xFF3F} & wave RAM & \cs{gb\_apu\_read}/\cs{write}, unmasked
  reads\\
\cs{0xFF40} & \cs{LCDC} & \cs{gb\_ppu\_read}/\cs{write}, power transitions\\
\cs{0xFF41} & \cs{STAT} & \cs{gb\_ppu\_read}/\cs{write}, bits 3--6 only\\
\cs{0xFF42} & \cs{SCY} & \cs{gb\_ppu\_read}/\cs{write}\\
\cs{0xFF43} & \cs{SCX} & \cs{gb\_ppu\_read}/\cs{write}\\
\cs{0xFF44} & \cs{LY} & \cs{gb\_ppu\_read}; writes ignored\\
\cs{0xFF45} & \cs{LYC} & \cs{gb\_ppu\_read}/\cs{write}\\
\cs{0xFF46} & \cs{DMA} & Special: starts an OAM transfer\\
\cs{0xFF47} & \cs{BGP} & \cs{gb\_ppu\_read}/\cs{write}\\
\cs{0xFF48} & \cs{OBP0} & \cs{gb\_ppu\_read}/\cs{write}\\
\cs{0xFF49} & \cs{OBP1} & \cs{gb\_ppu\_read}/\cs{write}\\
\cs{0xFF4A} & \cs{WY} & \cs{gb\_ppu\_read}/\cs{write}\\
\cs{0xFF4B} & \cs{WX} & \cs{gb\_ppu\_read}/\cs{write}\\
\cs{0xFF4D} & \cs{KEY1} & Special: bit~7 from \cs{double\_speed}, bit~0
  stored\\
\cs{0xFFFF} & \cs{IE} & \cs{bus->ie}, all eight bits\\
\bottomrule
\end{longtable}

Memory regions that are decoded but not emulated in full:

\begin{longtable}{@{}C{0.24\textwidth}L{0.66\textwidth}@{}}
\toprule
\textbf{Range} & \textbf{State}\\
\midrule
\endhead
\cs{0xFE00--0xFEA0} & OAM; reads return \cs{0xFF} and writes are dropped
  while a DMA transfer is copying\\
\cs{0xFEA0--0xFEFF} & Unusable; reads \cs{0xFF}, writes dropped\\
\cs{0xFF68}, \cs{0xFF69} & CGB palette; not decoded, lands in the shadow
  array\\
\cs{0xFF51--0xFF55} & CGB DMA; not decoded\\
\bottomrule
\end{longtable}

\clearpage

% ===========================================================================
\section{Public API reference}
\label{app-api}

The complete surface of \fname{src/gbemu.h}, in declaration order.

\subsection{CPU}

\begin{longtable}{@{}L{0.34\textwidth}L{0.58\textwidth}@{}}
\toprule
\textbf{Function} & \textbf{Purpose}\\
\midrule
\endhead
\cs{gb\_cpu\_init} & Set \cs{PC} to \cs{0x0100} and the post-boot
  registers\\
\cs{gb\_cpu\_step} & Fetch, decode and execute one instruction; return its
  T-cycles, or 4 while halted\\
\cs{gb\_cpu\_dump\_fd} & Print the register file to a \cs{FILE*}\\
\cs{gb\_flag\_set} / \cs{gb\_flag\_get} & Set or test one flag\\
\cs{gb\_get\_reg\_by\_index} / \cs{gb\_set\_reg\_by\_index} &
  Operand access by the 3-bit index, including the \cs{(HL)} machine
  cycle\\
\cs{gb\_fetch8} / \cs{gb\_fetch16} & Operand fetch, one machine cycle per
  byte\\
\cs{gb\_instr} & Decode and dispatch one opcode; returns its T-cycles\\
\bottomrule
\end{longtable}

\subsection{Bus}

\begin{longtable}{@{}L{0.34\textwidth}L{0.58\textwidth}@{}}
\toprule
\textbf{Function} & \textbf{Purpose}\\
\midrule
\endhead
\cs{gb\_bus\_read} / \cs{gb\_bus\_write} & The whole memory map\\
\cs{gb\_bus\_tick} & Advance the OAM DMA by one machine cycle\\
\cs{gb\_bus\_load\_rom} & Load and pad a ROM, size the SRAM, configure the
  mapper, sanitise the interrupt vectors; returns the byte count or 0\\
\cs{gb\_machine\_tick} & Advance DMA, timer, PPU and APU\\
\bottomrule
\end{longtable}

\subsection{Timer, PPU and APU}

\begin{longtable}{@{}L{0.34\textwidth}L{0.58\textwidth}@{}}
\toprule
\textbf{Function} & \textbf{Purpose}\\
\midrule
\endhead
\cs{gb\_timer\_init} / \cs{gb\_timer\_step} & Initialise and advance the
  divider and timer\\
\cs{gb\_timer\_read} / \cs{gb\_timer\_write} & The four registers\\
\cs{gb\_get\_time\_ns} / \cs{gb\_sleep\_ns} & Portable monotonic clock and
  sleep\\
\cs{gb\_ppu\_init} / \cs{gb\_ppu\_step} & Initialise and advance the
  scanline state machine\\
\cs{gb\_ppu\_read} / \cs{gb\_ppu\_write} & The video registers\\
\cs{gb\_apu\_init} / \cs{gb\_apu\_step} & Initialise and advance the
  synthesis\\
\cs{gb\_apu\_read} / \cs{gb\_apu\_write} & The sound registers\\
\cs{gb\_apu\_buf\_pop} & The one function a host may call from another
  thread\\
\cs{gb\_handle\_interrupts} & Latch request flags into \cs{IF}, wake the
  CPU, dispatch; returns 0 or 20\\
\bottomrule
\end{longtable}

\subsection{Loop and persistence}

\begin{longtable}{@{}L{0.34\textwidth}L{0.58\textwidth}@{}}
\toprule
\textbf{Function} & \textbf{Purpose}\\
\midrule
\endhead
\cs{gb\_loop} & The frame-budgeted main loop\\
\cs{gb\_battery\_load} / \cs{gb\_battery\_save} & The \cs{.sav} file\\
\cs{gb\_save\_state} / \cs{gb\_load\_state} & The \cs{.state} file\\
\bottomrule
\end{longtable}

\seealso
\fname{src/frontend.h} adds the \cs{gb\_frontend} struct, the
\cs{gb\_hotkey} enum and \cs{gb\_version}, which expand to the version
string, \cs{0.1.0}, and \cs{GB\_VERSION\_HEX}, which expands to \cs{0x100}.

\clearpage

% ===========================================================================
\section{Constants}
\label{app-constants}

Every compile-time constant in the public header, with the value it has
in this version.

\begin{longtable}{@{}L{0.38\textwidth}C{0.20\textwidth}L{0.34\textwidth}@{}}
\toprule
\textbf{Macro} & \textbf{Value} & \textbf{Meaning}\\
\midrule
\endhead
\cs{GB\_CPU\_FREQ} & \cs{4194304} & T-cycles per second\\
\cs{GB\_CYCLES\_PER\_FRAME} & \cs{69905} & \cs{GB\_CPU\_FREQ / 60}\\
\cs{GB\_FRAME\_TIME\_NS} & \cs{16666666} & one second, divided by 60\\
\cs{GB\_SCREEN\_WIDTH} & \cs{160} & Frame buffer width\\
\cs{GB\_SCREEN\_HEIGHT} & \cs{144} & Frame buffer height\\
\cs{GB\_APU\_BUF\_SIZE} & \cs{4096} & Ring buffer capacity, a power of two\\
\cs{GB\_APU\_BUF\_MASK} & \cs{4095} & \cs{SIZE - 1}, for wrapping\\
\cs{GB\_APU\_SAMPLE\_RATE} & \cs{44100} & Output sample rate\\
\cs{GB\_VERSION\_MAJOR} & \cs{0} &\\
\cs{GB\_VERSION\_MINOR} & \cs{1} &\\
\cs{GB\_VERSION\_PATCH} & \cs{0} &\\
\cs{GB\_VERSION\_HEX} & \cs{0x100} & \cs{0x00MMmmpp}, written into save
  states\\
\bottomrule
\end{longtable}

PPU modes and CPU flags:

\begin{longtable}{@{}L{0.38\textwidth}C{0.20\textwidth}L{0.34\textwidth}@{}}
\toprule
\textbf{Constant} & \textbf{Value} & \textbf{Meaning}\\
\midrule
\endhead
\cs{GB\_PPU\_MODE\_HBLANK} & \cs{0} & Mode~0, \cs{STAT} bits~1--0
  \cs{00}\\
\cs{GB\_PPU\_MODE\_VBLANK} & \cs{1} & Mode~1, \cs{01}\\
\cs{GB\_PPU\_MODE\_OAM} & \cs{2} & Mode~2, \cs{10}\\
\cs{GB\_PPU\_MODE\_XFER} & \cs{3} & Mode~3, \cs{11}\\
\cs{GB\_FLAG\_Z} & \cs{0x80} & Bit~7 of \fname{F}\\
\cs{GB\_FLAG\_N} & \cs{0x40} & Bit~6\\
\cs{GB\_FLAG\_H} & \cs{0x20} & Bit~5\\
\cs{GB\_FLAG\_C} & \cs{0x10} & Bit~4\\
\bottomrule
\end{longtable}

\clearpage

% ===========================================================================
\section{Further reading}
\label{app-reading}

The specifications and test ROMs this emulator is written against.  All
of them are free and public.

\subsection{Reference documentation}

\begin{longtable}{@{}L{0.30\textwidth}L{0.62\textwidth}@{}}
\toprule
\textbf{Source} & \textbf{Covers}\\
\midrule
\endhead
\cs{gbdev.io/pandocs/PPU.html} &
  Mode timing, \cs{STAT}, \cs{LY}/\cs{LYC}, the background, window and
  sprite rules, and the CGB additions\\
\cs{gbdev.io/pandocs/Memory\_Map.html} &
  The address map, the I/O register semantics and the DMA rules\\
\cs{gbdev.io/pandocs/Timer\_and\_Divider\_Registers.html} &
  The divider bit select, the \cs{TIMA} overflow sequence and the
  \cs{DIV}-write side effects\\
\cs{gbdev.io/pandocs/Audio.html} &
  The four channels, the frame sequencer steps, the sweep overflow rules
  and the length-counter timing\\
\cs{gbdev.io/pandocs/Cartridge\_Header.html} &
  Header bytes \cs{0x0143}, \cs{0x0147}, \cs{0x0148}, \cs{0x0149} and
  the Nintendo logo\\
\cs{gbdev.io/pandocs/Interrupt\_Handling.html} &
  The \cs{IF}/\cs{IE} contract, the five vectors and the five machine
  cycles of dispatch\\
\cs{gbdev.io/pandocs/MBC1.html}, \cs{MBC2}, \cs{MBC3}, \cs{MBC5} &
  One page per implemented mapper\\
\bottomrule
\end{longtable}

\subsection{Test ROMs}

\begin{longtable}{@{}L{0.30\textwidth}L{0.62\textwidth}@{}}
\toprule
\textbf{Suite} & \textbf{Note}\\
\midrule
\endhead
\cs{gb-test-roms} &
  Blargg's ROMs, vendored under \cs{tests/}.  The most widely used
  functional suite; the source of most of the known gaps in
  section~\ref{known-gaps}.\\
\cs{SameSuite} &
  SameBoy's suite, vendored under \cs{tests/}.  One ROM per hardware
  behaviour, which makes it the best tool for localising a bug.\\
\cs{mooneye-test-suite} &
  The acceptance tests, which are the closest thing to a specification
  that can be executed.  Several tests only pass on hardware; the harness
  filters by the suffix in the name.\\
\bottomrule
\end{longtable}

\subsection{Internal notes}

Within the repository, \fname{docs/dev/} holds the developer-facing
material that is not part of this manual: the header-comment markers for
\cs{tinydocs-cli}, the opcode JSON the generator consumes, and progress
notes.  \cs{make docs} regenerates the API pages from the headers.

\clearpage

% ===========================================================================
\section{Index}

\printindex

\end{document}
